\documentclass[11pt,a4paper]{article}
\usepackage[margin=2.6cm]{geometry}
\usepackage{amsmath,amssymb,amsthm}
\usepackage{booktabs}
\usepackage{array}
\usepackage{enumitem}
\usepackage{xcolor}
\usepackage[hidelinks]{hyperref}
\usepackage{microtype}
\usepackage{tabularx}
\setlength{\emergencystretch}{3em}

\newcommand{\G}{\mathbb{G}}
\newcommand{\Z}{\mathbb{Z}}
\newcommand{\Enc}{\mathsf{Enc}}
\newcommand{\Dec}{\mathsf{Dec}}
\newcommand{\com}{\mathsf{com}}
\newcommand{\pos}{\mathit{pos}}
\newcommand{\rocq}[1]{\texttt{#1}}
\newcommand{\ex}[1]{\par\medskip\noindent\textbf{Example.}~#1\par\medskip}
\newtheorem{theorem}{Theorem}[section]
\newtheorem{lemma}[theorem]{Lemma}
\newtheorem{definition}[theorem]{Definition}
\newtheorem{corollary}[theorem]{Corollary}
\newtheorem{proposition}[theorem]{Proposition}
\theoremstyle{remark}
\newtheorem*{remark}{Remark}

\title{Two Certified Schulze Ballots, Explained\\[4pt]
\large A tutorial on the receipt-free Schulze election of the paper, with worked examples and proofs}
\author{Companion note to the paper\\ \emph{A Publicly Tally-Hiding Schulze Election}\\ \emph{with Machine-Checked Soundness}}
\date{October 2026}

\begin{document}
\maketitle

\begin{abstract}
This note explains, from first principles and with numbers small enough to
check by hand, the two ballot encodings and the tally certificate of the
paper's Schulze election: how a voter encrypts a ranking and proves that the
ciphertexts are those of one ranking, how the pairwise margins are
aggregated without a key, how the trustees certify the Schulze winner over
encrypted margins without publishing them, and why the whole thing is sound,
zero-knowledge and receipt-free. It covers the cryptography the scheme rests
on (exponential ElGamal, Pedersen vector commitments, $\Sigma$-protocols and
their OR-composition, the Fiat--Shamir transform, a $\Sigma$-protocol for the
preimage of a linear map, and the Bayer--Groth zero argument), then walks
through the first ballot (membership proofs, $O(m^3)$ cost), the batched
ballot (four vector arguments, $O(m^2 \log m)$ cost), and the closure
certificate, each on a three-candidate election with five voters in a toy
group of order $23$, with every proof written out. It ends with the list of
machine-checked theorems and what each one says. It is written for readers
who know what a public-key cryptosystem is and have not worked with
zero-knowledge proofs.
\end{abstract}

\tableofcontents

\section{How to read this note}\label{sec:howto}
The paper states the scheme compactly and relies on the Rocq development for
the proofs. This note is the long version: it takes one election, three
candidates $A, B, C$ and five voters, and follows it through every step of
both protocols, with the group small enough that each equation can be
recomputed with a pocket calculator. The proofs are given in full where the
paper gives a sketch. Rocq names are in \rocq{typewriter}; every theorem
quoted here is proved in the development and listed in its
\rocq{THEOREMS.md}, and Section~\ref{sec:mech} maps the statements of this
note to those names.

The order is: the problem (Section~\ref{sec:problem}); the cryptographic
tools (Section~\ref{sec:tools}); the three ballots step by step, with
the toy numbers but without proofs (Section~\ref{sec:steps}); the first ballot
(Section~\ref{sec:ballot1}); aggregation (Section~\ref{sec:agg}); the
batched ballot, with a proposed simplification at its end
(Section~\ref{sec:ballot2}); one ballot through both protocols with every
number (Section~\ref{sec:walk}); the tally certificate
(Section~\ref{sec:tally}); security (Section~\ref{sec:security}); a
general recipe for designing such proofs (Section~\ref{sec:recipe});
where the ideas come from (Section~\ref{sec:sources}); the mechanisation
(Section~\ref{sec:mech}); and running the artefact
(Section~\ref{sec:artefact}). A reader who wants only the second protocol
can read Sections~\ref{sec:problem}, \ref{sec:tools} and~\ref{sec:ballot2};
a reader who wants the protocols without the proofs can read
Sections~\ref{sec:tools} and~\ref{sec:steps}.

\paragraph{The toy group} Throughout, the examples use the Schnorr group of
order $q = 23$ inside $\Z_{47}^*$: $p = 47 = 2 \cdot 23 + 1$, and $g = 2$
has order $23$ modulo $47$ (it is a quadratic residue, since $47 \equiv -1
\pmod 8$, and $2 \ne 1$). All exponents live in $\Z_{23}$ and all group
elements are integers modulo $47$. The election key is $x = 5$, $h = g^5 =
32$. Real elections use a $2048$-bit $p$ with a $256$-bit $q$ (the Helios
parameters) and nothing else changes.

\paragraph{The running example} Three candidates $A, B, C$, five voters:
\begin{gather*}
  \text{ballot 1: } A > B = C,\qquad
  \text{ballots 2, 3: } B > C > A,\\
  \text{ballot 4: } C > A > B,\qquad
  \text{ballot 5: } A > C > B.
\end{gather*}
We will encode ballot 1 in both protocols, aggregate all five, and certify
the winner, which turns out to be $C$.

\section{The problem}\label{sec:problem}

\subsection{Ranked ballots and the Schulze method}
A ranked ballot over $m$ candidates orders them, with ties allowed; Gentoo
and Debian also allow leaving candidates unranked, which counts as ranking
them last. Its pairwise content is, for each ordered pair $(i, j)$, whether
the voter ranks $i$ strictly above $j$. Summed over all ballots this gives
the \emph{margin matrix} $D$: $D_{ij}$ is the number of voters who rank $i$
above $j$. Everything the Schulze method does depends on the ballots only
through $D$.

From $D$, candidate $i$ \emph{defeats} $j$ with strength $d_{ij} = D_{ij}$
if $D_{ij} > D_{ji}$, and $d_{ij} = 0$ otherwise. A \emph{path} from $i$ to
$j$ is a sequence $i = u_0, u_1, \ldots, u_\ell = j$ with $\ell \ge 1$; its
strength is its weakest defeat, $\min_t d_{u_t u_{t+1}}$. The
\emph{strongest-path} matrix $p$ has $p_{ij}$ equal to the maximum strength
of a path from $i$ to $j$ (and $0$ if no path has positive strength).
Candidate $w$ is a \emph{Schulze winner} if $p_{wj} \ge p_{jw}$ for every
$j$. At least one winner always exists; with ties among winners, Gentoo and
Debian break them by a published rule, which we do not touch.

\ex{The five ballots give the margins (rows are $i$, columns $j$):
\[
  D = \begin{pmatrix} - & 3 & 2 \\ 2 & - & 2 \\ 3 & 2 & - \end{pmatrix},
\]
for instance $D_{AB} = 3$ because ballots $1, 4, 5$ rank $A$ above $B$, and
$D_{BC} = D_{CB} = 2$: ballots $2, 3$ have $B > C$, ballots $4, 5$ have $C >
B$, ballot $1$ ties them. The defeats are $d_{AB} = 3$ (since $3 > 2$) and
$d_{CA} = 3$ (since $3 > 2$); every other pair has $D_{ij} \le D_{ji}$, so
$d_{ij} = 0$. The strongest paths: $p_{AB} = 3$ (the edge), $p_{CA} = 3$
(the edge), $p_{CB} = \min(d_{CA}, d_{AB}) = 3$ via $C \to A \to B$, and
$p_{AC} = p_{BA} = p_{BC} = 0$, since every path from $A$ to $C$, from $B$
to $A$ or from $B$ to $C$ uses an edge of strength $0$. The winner is $C$:
$p_{CA} = 3 \ge p_{AC} = 0$ and $p_{CB} = 3 \ge p_{BC} = 0$. Candidate $A$
is not a winner, since $p_{AC} = 0 < p_{CA} = 3$.}

\subsection{What a published tally leaks}
A verifiable election must let anyone check the result. The simplest way is
to publish the ballots, as Gentoo does; but with $m! \cdot$(ties) possible
rankings, almost every ballot is unique, and a coercer who prescribes a
ranking can find it on the board (the ``Italian attack''). Publishing only
$D$ is better but, in an electorate of a hundred, a margin of one or two
confirms a prescribed preference between two candidates. The paper measures
this on Gentoo's thirteen council elections: in every one, all but at most
two ballots are unique rankings, and up to ten candidate pairs have a margin
of at most two.

\subsection{What we want}
We want a scheme in which: anyone can verify the result from the public
board with no key (\emph{universal verifiability}); the public learns the
winner and nothing else about the ballots (\emph{public tally-hiding}, in
the terminology of Kryvos); the trustees, who hold the decryption key,
learn the margin matrix $D$ but no ballot; and nothing on the board lets a
voter prove to a coercer how they voted (\emph{receipt-freeness} against a
coercer who sees the whole board). We also want the costly parts to be
cheap enough to run on real elections, and the soundness of the two checks
that decide the result, the ballot check and the certificate check, to be
proved in a proof assistant.

\section{Cryptographic tools}\label{sec:tools}

\subsection{Exponential ElGamal}
Fix the group $\G = \langle g \rangle$ of prime order $q$ and the key $h =
g^x$. The \emph{exponential ElGamal} encryption of a small integer $v$ with
randomness $r \in \Z_q$~\cite{elgamal} is
\[
  \Enc(v; r) = (g^r,\; g^v h^r).
\]
Decryption with $x$ computes $c_2 \cdot c_1^{-x} = g^v h^r g^{-rx} = g^v$,
and recovers $v$ from $g^v$ by trying $v = 0, 1, 2, \ldots$ up to a known
bound; in this scheme every plaintext is a bit or a count of voters, so the
bound is the number of voters $N$. We write $\Dec(c) = g^v$ for the group
element, which is what the proofs talk about.

The encryption is \emph{homomorphic}: multiplying ciphertexts componentwise
adds plaintexts and randomness,
\[
  \Enc(v; r) \cdot \Enc(v'; r') = \Enc(v + v'; r + r'),
  \qquad
  \Enc(v; r)^a = \Enc(a v; a r),
\]
and dividing subtracts. The security property we need is IND-CPA under the
decisional Diffie--Hellman assumption: without $x$, ciphertexts of different
plaintexts are indistinguishable.

\ex{With $g = 2$, $h = 32$ in $\Z_{47}$: $\Enc(1; 3) = (2^3, 2^1 \cdot
32^3) = (8, 18)$ and $\Enc(0; 7) = (2^7, 32^7) = (34, 7)$. Decrypting the
first: $18 \cdot 8^{-5} = 2 = g^1$. Their product is $(8 \cdot 34, 18
\cdot 7) = (37, 32) = \Enc(1; 10)$, and their quotient is $(8 \cdot
34^{-1}, 18 \cdot 7^{-1}) = (3, 16) = \Enc(1; 3 - 7) = \Enc(1; 19)$.}

\subsection{Pedersen vector commitments}
Fix a second generator $h_1$ and generators $g_1, \ldots, g_n$ of $\G$,
chosen so that nobody knows discrete logarithms among them (in practice:
hashed to the group from public labels). The \emph{Pedersen vector
commitment}~\cite{pedersen} to $v \in \Z_q^n$ with randomness $\rho$ is
\[
  \com(v; \rho) = h_1^{\rho} \prod_{j=1}^{n} g_j^{v_j}.
\]
It is \emph{perfectly hiding}: for uniform $\rho$ the commitment is uniform
in $\G$ whatever $v$ is. It is \emph{binding} under the discrete-logarithm
assumption: two openings $(v, \rho) \ne (v', \rho')$ of the same commitment
give $h_1^{\rho - \rho'} \prod_j g_j^{v_j - v'_j} = 1$, a non-trivial
relation among the generators, which nobody should be able to find. The
theorems of the development conclude ``\ldots or a binding break'', i.e.\
they exhibit such a relation when extraction fails
(\rocq{binding\_break}). It is \emph{homomorphic}: $\com(v; \rho) \cdot
\com(v'; \rho') = \com(v + v'; \rho + \rho')$, so a verifier can derive a
commitment to $v - \mathbf{1}$ from one to $v$ by multiplying with
$\com(-\mathbf{1}; 0)$.

\ex{Take $h_1 = 32$, $g_1 = 34$, $g_2 = 8$ (all in the order-$23$
subgroup). Then $\com((1, 0); 4) = 32^4 \cdot 34 = 16$, $\com((0, 1); 2) =
32^2 \cdot 8 = 14$, and their product $16 \cdot 14 = 36 = \com((1, 1); 6)$.}

\subsection{$\Sigma$-protocols}\label{sec:sigma}
A \emph{$\Sigma$-protocol} for a relation $R$ (statement $T$, witness $w$)
is a three-move interaction: the prover sends an \emph{announcement} $A$,
the verifier a uniformly random \emph{challenge} $c \in \Z_q$, the prover a
\emph{response} $z$, and the verifier accepts or rejects $(T, A, c, z)$.
Three properties:
\begin{itemize}[nosep]
  \item \emph{Completeness}: an honest prover with a witness is accepted.
  \item \emph{Special soundness}: from two accepting transcripts $(A, c,
  z)$ and $(A, c', z')$ with the same announcement and $c \ne c'$, a
  witness can be computed. This is what the mechanised soundness theorems
  state: they take the two transcripts as hypotheses and output the
  witness (``rewinding form''). The on-paper step that turns it into ``a
  prover accepted with probability $\varepsilon$ knows a witness'' is the
  forking lemma.
  \item \emph{Special honest-verifier zero-knowledge (SHVZK)}: for every
  challenge $c$ there is a simulator that, from $T$ and $c$ alone, outputs
  $(A, z)$ distributed exactly as in a real transcript with that
  challenge. The development states this as a \emph{permutation}: the list
  of (transcript, probability) pairs of the real prover over its
  randomness is a permutation of the simulator's list. Perfect, not
  computational.
\end{itemize}

\paragraph{Schnorr} Statement $T = g^w$; announcement $A = g^a$ for random
$a$; response $z = a + c w$; check $g^z = A \cdot T^c$. Completeness:
$g^{a + cw} = g^a (g^w)^c$. Special soundness: from $g^z = A T^c$ and
$g^{z'} = A T^{c'}$, dividing, $g^{z - z'} = T^{c - c'}$, so $w = (z -
z')/(c - c')$. SHVZK: the simulator picks $z$ uniformly and sets $A = g^z
T^{-c}$; the map $a \mapsto z = a + cw$ is a bijection of $\Z_q$, so real and
simulated transcripts have the same distribution.

\ex{$w = 5$, $T = h = 32$. Prover: $a = 9$, $A = 2^9 = 42$. Challenge $c =
6$. Response $z = 9 + 6 \cdot 5 = 39 \equiv 16 \pmod{23}$. Check: $2^{16} =
18$ and $42 \cdot 32^6 = 18$.}

\paragraph{Chaum--Pedersen} The same with two bases: prove that $(T_1, T_2)
= (g^w, k^w)$ for public $g, k$: $A = (g^a, k^a)$, $z = a + cw$, check
$g^z = A_1 T_1^c$ and $k^z = A_2 T_2^c$. This proves that a ciphertext
$(c_1, c_2)$ encrypts a given $v$: $(c_1, c_2 / g^v) = (g^r, h^r)$ is a
Chaum--Pedersen statement with witness $r$ and bases $g, h$. It also proves
correct decryption: $(h, \Dec(c)^{-1} c_2) = (g^x, c_1^x)$ with witness
$x$. The development calls these \rocq{EncProof} (for one message) and
\rocq{DecProof}.

\paragraph{OR-composition and membership proofs} To prove that a ciphertext
encrypts \emph{one of} $v_1, \ldots, v_n$ without saying which, the prover
runs the Chaum--Pedersen protocol for the true branch and \emph{simulates}
the others (Cramer, Damg{\aa}rd, and Schoenmakers~\cite{cds}): it picks challenges
$c_k$ and responses $z_k$ for every false branch $k$ and computes their
announcements from the verification equation; it computes the true branch's
announcement honestly; after receiving the challenge $c$ it sets the true
branch's challenge to $c - \sum_{k \ne \text{true}} c_k$ and answers it.
The verifier checks every branch and that the branch challenges sum to
$c$. Special soundness: two transcripts with $c \ne c'$ differ in at least
one branch's challenge, and that branch yields a witness, i.e.\ the
randomness $r$ of $\Enc(v_k; r)$ for some $k$; so the ciphertext encrypts
one of the messages. SHVZK: simulate all branches. The cost is linear in
$n$: one Chaum--Pedersen transcript per branch. This is \rocq{EncProof}
over a list of messages, and ``the plaintext is in a public set'' is called
a \emph{membership proof} below.

\ex{Prove that $C = \Enc(1; 3) = (8, 18)$ encrypts $0$ or $1$. Branch $0$
(false) is simulated: pick $c_0 = 4$, $z_0 = 11$, set $A_0 = (g^{z_0}
C_1^{-c_0}, h^{z_0} C_2^{-c_0}) = (24, 28)$. Branch $1$ (true) is honest:
$a_1 = 10$, $A_1 = (g^{10}, h^{10}) = (37, 16)$. Challenge $c = 9$; then
$c_1 = 9 - 4 = 5$ and $z_1 = 10 + 5 \cdot 3 = 25 \equiv 2$. The verifier
checks $c_0 + c_1 = 9$, branch $0$: $g^{11} = A_{0,1} C_1^{4}$ and
$h^{11} = A_{0,2} C_2^{4}$ (true by construction), branch $1$: $g^{2} = 4 =
37 \cdot 8^{5}$ and $h^{2} = 37 = 16 \cdot (18 / 2)^{5}$. Both hold.}

\paragraph{Fiat--Shamir} The interaction is removed by computing the
challenge as a hash of the statement and the announcement, $c = H(T,
A)$~\cite{fiatshamir};
the verifier recomputes it. In the ballot the hash also takes the voter's
credential, so that a ballot copied under another credential is rejected.
The mechanised theorems are stated for an arbitrary function $H$; the
random-oracle argument that turns special soundness into non-interactive
soundness is on paper.

\subsection{A $\Sigma$-protocol for the preimage of a linear map}\label{sec:linsigma}
Let $V$ be a vector space over $\Z_q$ (written additively) and $\varphi :
\Z_q^W \to V$ a linear map. To prove knowledge of $w$ with $\varphi(w) =
T$, the protocol of Maurer~\cite{maurer} is:
\[
  A = \varphi(a) \text{ for random } a \in \Z_q^W, \qquad
  z = a + c\,w, \qquad
  \text{check } \varphi(z) = A + c\,T .
\]
Schnorr is the case $V = \G$ (written multiplicatively), $\varphi(w) =
g^w$; Okamoto's representation proof is $\varphi(w_1, w_2) = g^{w_1}
h^{w_2}$; and a product $V = \G^{k} \times \Z_q^{l}$ lets one proof cover
$k$ group equations and $l$ field equations with a single challenge and a
single response vector. This is \rocq{Crypto/LinearSigma.v}.

\begin{lemma}[\rocq{lin\_completeness}, \rocq{lin\_soundness}, \rocq{lin\_shvzk}]
The protocol is complete, special sound with extractor $w = (z - z')/(c -
c')$, and perfectly SHVZK.
\end{lemma}
\begin{proof}
Completeness: $\varphi(a + cw) = \varphi(a) + c\,\varphi(w) = A + cT$ by
linearity. Special soundness: subtracting the two verification equations,
$\varphi(z - z') = (c - c') T$, so $\varphi\big((z - z')/(c - c')\big) =
T$ by linearity, since $c - c' \ne 0$ is invertible in $\Z_q$. SHVZK: the
simulator draws $z$ uniformly in $\Z_q^W$ and sets $A = \varphi(z) - cT$;
in the real protocol $z = a + cw$ with $a$ uniform, and $a \mapsto a + cw$
is a bijection of $\Z_q^W$ with inverse $z \mapsto z - cw$, under which the
real announcement $\varphi(a) = \varphi(z - cw) = \varphi(z) - cT$ is the
simulated one. So the two distributions are the same list up to the order
in which the randomness is enumerated (\rocq{zk\_of\_bijection}).
\end{proof}

\subsection{The Bayer--Groth zero argument}\label{sec:zero}
For a challenge $y \in \Z_q$ define the bilinear map on vectors of length
$n$,
\[
  a \star_y b = \sum_{j=1}^{n} y^{j} a_j b_j .
\]
The \emph{zero argument} takes commitments $C_{a_1}, \ldots, C_{a_r}$ and
$C_{b_1}, \ldots, C_{b_r}$ to vectors $a_i, b_i \in \Z_q^n$ and proves
$\sum_{i=1}^{r} a_i \star_y b_i = 0$. It is the sub-argument of the
Hadamard argument in Bayer and Groth's shuffle~\cite{bayergroth}, built on
the bilinear-map technique of Groth~\cite{groth09}; the development mechanises
it in the form of the Swiss Post specification (\rocq{BayerGroth.v},
\rocq{zero\_verify}, \rocq{zero\_prove}). We use it as a black box with
three properties:
\begin{itemize}[nosep]
  \item completeness for honest openings satisfying the relation
  (\rocq{zero\_completeness});
  \item $(2r{+}1)$-special soundness: from $2r+1$ accepting transcripts
  with distinct challenges $x$ on the same announcement, openings of all
  $C_{a_i}, C_{b_i}$ satisfying $\sum_i a_i \star_y b_i = 0$ are
  extracted, or a binding break of the commitment key
  (\rocq{zero\_soundness}, \rocq{BayerGrothSound.v});
  \item perfect SHVZK, given the discrete logarithms of the commitment key
  in base $h_1$ (\rocq{zero\_shvzk}).
\end{itemize}
Why the relation is useful: read $\sum_i a_i \star_y b_i = \sum_j y^j
\big(\sum_i a_{ij} b_{ij}\big)$ as a polynomial in $y$ of degree $n$ whose
coefficient at $y^j$ is $\sum_i a_{ij} b_{ij}$. A non-zero polynomial of
degree $n$ has at most $n$ roots in the field (the one-variable case of
the Schwartz--Zippel lemma~\cite{schwartz,zippel}). So if the relation holds for
$n$ distinct non-zero values of $y$, every coefficient is zero, i.e.\
$\sum_i a_{ij} b_{ij} = 0$ \emph{for every coordinate $j$}. With a single
random $y$ (Fiat--Shamir), the relation holds for a false statement with
probability at most $n/q$. Two instances matter for us:
\begin{itemize}[nosep]
  \item $r = 2$, $a = (b, p)$, $b = (\delta, -\mathbf{1})$: the coordinate
  relations are $b_j \delta_j - p_j = 0$, i.e.\ $p = b \circ \delta$
  (entry-wise product, ``Hadamard'');
  \item $r = 1$, $a = X$, $b = X - \mathbf{1}$: the relations are $X_j (X_j
  - 1) = 0$, i.e.\ every $X_j$ is a bit.
\end{itemize}

The argument itself is short. Take one row, $r = 1$: the claim is $a
\star_y c = 0$ for committed vectors $a$ and $c$. The prover draws
blinding vectors $a_0, c_0 \in \Z_q^n$ and sets
\[
  A(x) = a_0 + x\,a, \qquad C(x) = x\,c + c_0, \qquad
  A(x) \star_y C(x) = d_0 + x\,d_1 + x^2\,(a \star_y c),
\]
with $d_0 = a_0 \star_y c_0$ and $d_1 = a_0 \star_y c + a \star_y c_0$. It
announces $\com(a_0)$, $\com(c_0)$, and commitments $D_0, D_1$ to the
scalars $d_0, d_1$ (over the first generator: $\com(d; t) = h_1^t
g_1^d$), receives the challenge $x$, and answers with $A(x)$, $C(x)$, and
the combined randomness. The verifier checks three equations:
\[
  \com(A(x)) = \com(a_0) \cdot C_a^{\,x}, \qquad
  \com(C(x)) = C_c^{\,x} \cdot \com(c_0), \qquad
  \com\big(A(x) \star_y C(x)\big) = D_0 \cdot D_1^{\,x}.
\]
There is no commitment for the coefficient of $x^2$, so the last equation
can hold for many $x$ only if that coefficient, $a \star_y c$, is zero:
three accepting challenges on one announcement determine the
degree-$2$ polynomial, which is the $3$-special soundness. With $r$ rows
the prover uses polynomials of degree $r$ in $x$ and commits to the $2r$
coefficients other than the one claimed to be zero, which is why $2r + 1$
challenges are needed. The blinding vectors make $A(x)$ and $C(x)$
uniform, which is where zero-knowledge comes from. The development
mechanises the Swiss Post form of the argument~\cite{swisspost}, which
arranges the same polynomial with different indexing.

\section{The protocols, step by step}\label{sec:steps}
This section runs through each protocol once, from what the voter knows
to what the verifier checks: the first ballot, the batched ballot, and
the simplified batched ballot, which is the batched ballot with its two
linear checks moved into the verifier and its three zero arguments merged
into one. Every step uses only the tools of
Section~\ref{sec:tools} and the steps before it. The later sections give
the proofs, the costs, and the Rocq names; this one gives the protocols.
The numbers are in the toy group ($p = 47$, $q = 23$, $g = 2$, $x = 5$, $h
= 32$). The first protocol is shown on ballot 1, $A > B = C$, and the two
batched ones on $A > B > C$, whose ciphertexts and commitments are those of
Section~\ref{sec:walk}.

\subsection{What every protocol must prove}
The voter's ranking is a position vector $s \in \{0, \ldots, m\}^m$: $s_i$
is the position of candidate $i$, a smaller position is better, and equal
positions are ties. The tally needs, for every ordered pair $(i, j)$, the
encrypted bit
\[
  b_{ij} = [\, s_i < s_j \,] = [\, \delta_{ij} \ge 1 \,], \qquad
  \delta_{ij} = s_j - s_i \in \{-m, \ldots, m\}.
\]
A ballot is valid when every posted bit ciphertext encrypts $[s_i < s_j]$
for one and the same $s$. A dishonest voter has two ways to break this.
The first is to encrypt a value that is not a bit, such as $5$, which adds
five to a margin. The second is to encrypt bits that come from no single
ranking, such as $b_{AB} = b_{BA} = 1$ or the cycle $b_{AB} = b_{BC} =
b_{CA} = 1$. Both protocols therefore prove three things: every $b_{ij}$
is a bit, all bits are anchored to one position vector, and each bit
agrees with the sign of its difference. The first protocol proves them
ciphertext by ciphertext with disjunctions; the batched protocol proves
them for whole vectors with arithmetic.

\subsection{The first protocol}

\subsubsection*{Step 1: encrypt the positions and the bits}
The voter posts $\Enc(s_i; r_i)$ for every candidate $i$ and $\Enc(b_{ij};
r_{ij})$ for every ordered pair. Only the bit ciphertexts are aggregated;
the position ciphertexts exist so that the bits can be checked against
them.

\ex{Ballot 1 is $s = (0, 1, 1)$. With $r_A = 3$ and $r_B = 7$, the voter
posts $\Enc(s_A) = \Enc(0; 3) = (8, 9)$ and $\Enc(s_B) = \Enc(1; 7) = (34,
14)$. The bit $b_{AB} = 1$ is encrypted with randomness $5$: $\Enc(1; 5) =
(2^5, 2 \cdot 32^5) = (32, 8)$.}

\subsubsection*{Step 2: the verifier derives every difference}
The voter never posts $\delta$. The verifier computes it with the
homomorphism,
\[
  \Enc(\delta_{ij}) = \Enc(s_j) / \Enc(s_i) = \Enc(s_j - s_i;\; r_j - r_i).
\]
This is the \emph{anchoring}: all $m^2$ differences come from the same $m$
position ciphertexts, so $\delta_{BA} = -\delta_{AB}$ and $\delta_{AB} +
\delta_{BC} = \delta_{AC}$ hold automatically, and no voter can pick
differences that no ranking has.

\ex{$\Enc(\delta_{AB}) = (34, 14) / (8, 9) = (16, 12) = \Enc(1; 7 - 3) =
\Enc(1; 4)$.}

\subsubsection*{Step 3: merge the two cases into one ciphertext}
The bit must satisfy $(b_{ij} = 1 \wedge \delta_{ij} \ge 1) \vee (b_{ij} =
0 \wedge \delta_{ij} \le 0)$. To check both halves with a single
ciphertext, the verifier computes, again by the homomorphism, the
\emph{shifted} ciphertext
\[
  \Enc(b_{ij})^{K} \cdot \Enc(\delta_{ij}) = \Enc(K b_{ij} + \delta_{ij};\;
  K r_{ij} + r_j - r_i), \qquad K = 2m + 1 .
\]
Its plaintext $t = K b_{ij} + \delta_{ij}$ is in $\{-m, \ldots, 0\}$
exactly when $b_{ij} = 0$ and $\delta_{ij} \le 0$, and in $\{K + 1,
\ldots, K + m\}$ exactly when $b_{ij} = 1$ and $\delta_{ij} \ge 1$. The
gap $\{1, \ldots, K\}$ between the two ranges catches every wrong pair:
$b = 1$ with $\delta \le 0$ gives $t \in \{K - m, \ldots, K\}$, and $b =
0$ with $\delta \ge 1$ gives $t \in \{1, \ldots, m\}$.

\ex{$m = 3$, so $K = 7$ and the allowed set is $S_{\mathrm{pair}} = \{0,
-1, -2, -3, 8, 9, 10\}$. For $(A, B)$: $(32, 8)^7 \cdot (16, 12) = (7, 12)
\cdot (16, 12) = (18, 3) = \Enc(7 + 1; 7 \cdot 5 + 4) = \Enc(8; 16)$, and
$8 \in S_{\mathrm{pair}}$. A voter who claims $b_{BA} = 1$, encrypted as
$\Enc(1; 5)$, gets the shifted ciphertext $\Enc(1; 5)^7 \cdot (8, 9) / (34,
14) = (21, 1)$, which encrypts $7 - 1 = 6 \notin S_{\mathrm{pair}}$.}

\subsubsection*{Step 4: everything is now a membership claim}
After Steps 2 and 3, validity is three kinds of claim, each of the form
``this ciphertext encrypts one of a public list of messages'':
\begin{itemize}[nosep]
  \item every $\Enc(s_i)$ encrypts one of $0, 1, \ldots, m$ ($m + 1$
  messages);
  \item every $\Enc(b_{ij})$ encrypts one of $0, 1$ ($2$ messages);
  \item every shifted ciphertext encrypts one of $S_{\mathrm{pair}}$ ($2m
  + 1$ messages).
\end{itemize}
The bit claim is needed separately because the gap argument of Step 3
assumes $b \in \{0, 1\}$; the position claim is needed because it assumes
$|\delta| \le m$. Steps 5 to 8 build the membership proof from
Chaum--Pedersen proofs.

\subsubsection*{Step 5: one message is one Chaum--Pedersen proof}
A ciphertext $(C_1, C_2)$ encrypts $\mu$ with randomness $r$ exactly when
\[
  C_1 = g^{r} \quad\text{and}\quad C_2 / g^{\mu} = h^{r},
\]
one exponent $r$ for two bases. The Chaum--Pedersen proof of
Section~\ref{sec:sigma} proves this. The prover picks $u$ and announces $(a_1,
a_2) = (g^{u}, h^{u})$, receives a challenge $e$, and answers $z = u + e
r$. The verifier checks $g^{z} = a_1 C_1^{e}$ and $h^{z} = a_2 (C_2 /
g^{\mu})^{e}$. The single response $z$ for both equations is what forces
the same $r$: if $\mu$ is the wrong message, then $C_2 / g^{\mu} = h^{r}
g^{\mu_{\mathrm{true}} - \mu}$ is not a power of $h$ with exponent $r$,
and no $z$ satisfies both equations for more than one $e$.

\subsubsection*{Step 6: simulate the messages you did not encrypt}
For any $\mu$, true or false, an accepting transcript can be produced
without $r$ by choosing the challenge first. Pick $e$ and $z$, then set
\[
  a_1 = g^{z} / C_1^{e}, \qquad a_2 = h^{z} / (C_2 / g^{\mu})^{e}.
\]
Both checks of Step 5 then hold by construction. This is the simulator
behind SHVZK, and it is also what lets the prover handle the messages that
are not the plaintext.

\subsubsection*{Step 7: give the true message the leftover challenge}
To prove ``one of $\mu_1, \ldots, \mu_\ell$'' when the plaintext is
$\mu_{i^*}$, the prover:
\begin{enumerate}[nosep]
  \item simulates every false branch $i \ne i^*$ as in Step 6, choosing
  $e_i$ and $z_i$ and computing its announcement;
  \item announces the true branch honestly, $(g^{u}, h^{u})$;
  \item receives the single challenge $e$;
  \item sets $e_{i^*} = e - \sum_{i \ne i^*} e_i$ and answers $z_{i^*} = u
  + e_{i^*} r$;
  \item sends all announcements, all $e_i$, and all $z_i$.
\end{enumerate}
The verifier checks that $\sum_i e_i = e$ and that every branch passes the
two equations of Step 5 with its own $e_i$. Real and simulated branches
are identically distributed, and the $e_i$ are uniform apart from their
sum, so the transcript does not reveal $i^*$.

\ex{The bit proof for $\Enc(b_{AB}) = (32, 8)$, messages $\{0, 1\}$, true
branch $\mu = 1$ with $r = 5$, challenge $e = 15$.
\begin{center}
\begin{tabular}{@{}lll@{}}
\toprule
 & branch $\mu = 0$ (simulated) & branch $\mu = 1$ (true) \\
\midrule
prover picks & $e_0 = 9$, $z_0 = 11$ & $u = 6$ \\
announcement & $(2^{11} \cdot 32^{-9},\; 32^{11} \cdot 8^{-9}) = (7, 32)$ & $(2^6, 32^6) = (17, 34)$ \\
challenge & $e_0 = 9$ & $e_1 = 15 - 9 = 6$ \\
response & $z_0 = 11$ & $z_1 = 6 + 6 \cdot 5 = 36 \equiv 13$ \\
\bottomrule
\end{tabular}
\end{center}
The verifier checks $9 + 6 = 15$. For branch $1$: $g^{13} = 14$ and $17
\cdot 32^{6} = 17 \cdot 34 \equiv 14$; $h^{13} = 3$ and $34 \cdot (8 /
2)^{6} = 34 \cdot 7 \equiv 3$. Branch $0$ holds by construction.}

\subsubsection*{Step 8: why a false claim cannot be proved}
All announcements are fixed before $e$ is known. In a false branch the
prover can answer only the one challenge it chose in advance. If every
branch is false, the sum check $\sum_i e_i = e$ then holds for a single
value of $e$, probability $1/q$. The extractor states the same fact
exactly: two accepting transcripts with the same announcements and $e \ne
e'$ must differ in some branch's challenge, and that branch gives $r =
(z_i - z'_i) / (e_i - e'_i)$, so its message is the plaintext.

\ex{Rerun the bit proof above with the same announcements and $e' = 4$.
The simulated branch keeps $e_0 = 9$, so $e'_1 = 4 - 9 \equiv 18$ and
$z'_1 = 6 + 18 \cdot 5 \equiv 4$. From the two runs, $r = (13 - 4)/(6 -
18) = 9 / 11 \equiv 9 \cdot 21 \equiv 5$, the randomness of $\Enc(b_{AB})$;
so branch $\mu = 1$ is the true one.}

\subsubsection*{Step 9: prove every claim of the ballot}
The voter runs Step 7 once for each claim of Step 4. Each has its own true
branch.

\ex{The position proof for $\Enc(s_A) = (8, 9)$, messages $\{0, 1, 2,
3\}$, true branch $0$ with $r_A = 3$, $e = 15$. The simulated branches use
$(e_i, z_i) = (5, 2), (17, 9), (20, 14)$ for $i = 1, 2, 3$, with
announcements $(37, 42), (16, 21), (1, 28)$. The true branch announces
$(g^{10}, h^{10}) = (37, 16)$, receives $e_0 = 15 - (5 + 17 + 20) \equiv
19$, and answers $z_0 = 10 + 19 \cdot 3 \equiv 21$. Check: $g^{21} = 12 =
37 \cdot 8^{19}$ and $h^{21} = 14 = 16 \cdot 9^{19}$.

The consistency proof for $(A, B)$ is on the shifted ciphertext $(18, 3) =
\Enc(8; 16)$, with seven branches in the order of the message list:
\begin{center}
\begin{tabular}{@{}lccccccc@{}}
\toprule
branch $\mu$ & $0$ & $-1$ & $-2$ & $-3$ & $\mathbf{8}$ & $9$ & $10$ \\
\midrule
announcement & $(8,28)$ & $(27,17)$ & $(2,27)$ & $(17,4)$ & $(7,28)$ & $(16,34)$ & $(1,4)$ \\
$e_\mu$ & $3$ & $8$ & $4$ & $13$ & $\mathbf{22}$ & $10$ & $1$ \\
$z_\mu$ & $5$ & $1$ & $19$ & $7$ & $\mathbf{19}$ & $3$ & $16$ \\
\bottomrule
\end{tabular}
\end{center}
The true branch $\mu = 8$ uses $u = 12$, gets $e_8 = 15 - (3 + 8 + 4 + 13
+ 10 + 1) \equiv 22$, and answers $z_8 = 12 + 22 \cdot 16 \equiv 19$.
Check: $g^{19} = 3 = 7 \cdot 18^{22}$ and $h^{19} = 8 = 28 \cdot (3 /
g^{8})^{22} = 28 \cdot 27^{22}$. All seven challenges sum to $15$.}

\subsubsection*{Step 10: one challenge for the whole ballot}
In the non-interactive ballot every proof uses the same challenge $e = H(\ldots)$, the
hash of all ciphertexts and all announcements of the ballot
(\rocq{ballot\_challenge}). The voter first computes every announcement:
the simulated branches need only their own $e_i, z_i$, and the true
branches need only $u$. The voter then hashes and answers each true
branch. Changing any ciphertext or announcement changes $e$, so the voter
cannot choose the challenge after the fact.

\subsubsection*{Step 11: what the verifier does}
The verifier (\rocq{nizk\_verify\_ballot}):
\begin{enumerate}[nosep]
  \item recomputes $e$ and checks that every proof is bound to it
  (\rocq{challenges\_are});
  \item for every candidate $i$, checks the $(m+1)$-branch proof on
  $\Enc(s_i)$;
  \item for every pair $(i, j)$, checks the $2$-branch proof on
  $\Enc(b_{ij})$, computes $\Enc(\delta_{ij})$ and the shifted ciphertext,
  and checks the $(2m+1)$-branch proof on it;
  \item if everything passes, adds the $\Enc(b_{ij})$ to the encrypted
  margins.
\end{enumerate}
Accepted means valid: the position proofs fix one $s$, the bit proofs make
each $b_{ij}$ a bit, and the consistency proofs with the gap of Step 3
give $b_{ij} = [s_i < s_j]$ (Theorem~\ref{thm:b1sound}). Ballot 1 carries
$3 \cdot 4 + 9 \cdot 2 + 9 \cdot 7 = 93$ Chaum--Pedersen transcripts; in
general about $2m^3$, which is $6.2$\,MB at $m = 17$.

\subsection{The batched protocol}
The first protocol pays $2m^3$ transcripts because every comparison is a
disjunction over a range. The batched protocol removes the disjunctions:
it rewrites ``$b_{ij} = [s_i < s_j]$'' as linear equations, entry-wise
products, and ``this entry is a bit'', and proves each kind for a whole
committed vector at once. Write $n = m^2$, index cell $(i, j)$ by $k = im
+ j$, and let $L = \lceil \log_2(m + 1) \rceil$.

\subsubsection*{Step 1: replace the disjunction by arithmetic}
The difference vector is a public linear function of the positions,
$\delta = Ms$, where row $(i, j)$ of the $n \times m$ matrix $M$ is $e_j -
e_i$ ($-1$ in column $i$, $+1$ in column $j$). $M$ depends only on $m$,
and the verifier computes it as part of the protocol.

Because $\delta_{ij}$ is an integer, the two cases of validity are
\begin{itemize}[nosep]
  \item $b_{ij} = 1$: $\delta_{ij} \ge 1$, that is, $\delta_{ij} - 1 \ge 0$;
  \item $b_{ij} = 0$: $\delta_{ij} \le 0$, that is, $-\delta_{ij} \ge 0$.
\end{itemize}
One expression selects the right left-hand side for either value of the
bit:
\[
  v = b(\delta - 1) + (1 - b)(-\delta) = -\delta + 2\, b \circ \delta - b .
\]
So validity is equivalent to ``$b$ is a bit and $v \ge 0$''. A field has
no order, so $v \ge 0$ is proved by writing $v$ in binary:
\[
  v = \sum_{t < L} 2^{t} X^{(t)}, \qquad X^{(t)} \in \{0, 1\}^{n}.
\]
$L$ bits suffice because $v_k \le m$. The product $b \circ \delta$
multiplies two secret vectors, which no linear proof can check. The voter
therefore commits to it as a vector of its own, $p = b \circ \delta$, and
proves the product separately. The statement is
\[
  \delta = M s, \qquad p = b \circ \delta, \qquad v = -\delta + 2p - b =
  \sum_{t < L} 2^{t} X^{(t)}, \qquad b, X^{(0)}, \ldots, X^{(L-1)} \in
  \{0, 1\}^{n}.
\]
The first and third are linear in the committed vectors; the second and
the bit conditions are entry-wise products.

\ex{$A > B > C$, $s = (0, 1, 2)$, $L = 2$:
\begin{center}
\begin{tabular}{@{}lrrrrrrrrr@{}}
\toprule
 & $AA$ & $AB$ & $AC$ & $BA$ & $BB$ & $BC$ & $CA$ & $CB$ & $CC$ \\
\midrule
$\delta$ & $0$ & $1$ & $2$ & $-1$ & $0$ & $1$ & $-2$ & $-1$ & $0$ \\
$b$ & $0$ & $1$ & $1$ & $0$ & $0$ & $1$ & $0$ & $0$ & $0$ \\
$p$ & $0$ & $1$ & $2$ & $0$ & $0$ & $1$ & $0$ & $0$ & $0$ \\
$v$ & $0$ & $0$ & $1$ & $1$ & $0$ & $0$ & $2$ & $1$ & $0$ \\
$X^{(0)}$ & $0$ & $0$ & $1$ & $1$ & $0$ & $0$ & $0$ & $1$ & $0$ \\
$X^{(1)}$ & $0$ & $0$ & $0$ & $0$ & $0$ & $0$ & $1$ & $0$ & $0$ \\
\bottomrule
\end{tabular}
\end{center}
Cell $AC$: $v = -2 + 2 \cdot 2 - 1 = 1 = 1 + 2 \cdot 0$. Cell $CA$: $v = 2
+ 0 - 0 = 2 = 0 + 2 \cdot 1$.}

\subsubsection*{Step 2: encrypt the bits, commit to everything else}
The voter posts the $n$ ciphertexts $c_k = \Enc(b_k; r_k)$, the only values
that are aggregated, and five Pedersen vector commitments
(Section~\ref{sec:tools}):
\begin{center}
\begin{tabular}{@{}llll@{}}
\toprule
commitment & vector & length & key \\
\midrule
$C_s = \com(s; \rho_s)$ & positions & $m$ & $\mathbf{g}_1$ \\
$C_\delta = \com(\delta; \rho_\delta)$ & differences & $n$ & $\mathbf{g}_1$ \\
$C_b = \com(b; \rho_b)$ & bits & $n$ & $\mathbf{g}_1$ \\
$C_p = \com(p; \rho_p)$ & products & $n$ & $\mathbf{g}_1$ \\
$C_X = \com(X; \rho_X)$ & binary digits & $Ln$ & $\mathbf{g}_2$ \\
\bottomrule
\end{tabular}
\end{center}
The vector $X = X^{(0)} \,\|\, X^{(1)} \,\|\, \cdots \,\|\, X^{(L-1)}$ stacks
the bit vectors one after another, so entry $tn + k$ is bit $t$ of $v_k$,
and the verifier rebuilds $v_k = \sum_t 2^{t} X_{tn + k}$. The positions
are committed, not encrypted: nothing ever decrypts them.

$X$ is longer than the other vectors, so it needs a longer key; the
development uses an independent key $\mathbf{g}_2$ of length $Ln$. A single
key of length $Ln$, with the shorter vectors committed under its first
$n$ generators, would do as well. What soundness needs is that no
discrete-logarithm relation among $h_1$ and the generators is known, and
that the commitments combined in one zero argument (Step 5) share a key,
since that argument works entry by entry.

\ex{For $A > B > C$ and $m = 3$: $X = (0,0,1,1,0,0,0,1,0 \,\|\,
0,0,0,0,0,0,1,0,0)$, of length $18$. Cell $CA$ ($k = 6$) reads entries $6$
and $15$: $v_{CA} = 0 + 2 \cdot 1 = 2$. With the keys and randomness of
Section~\ref{sec:walk}, $C_s = 42$, $C_\delta = 27$, $C_b = 4$, $C_p = 17$,
and $C_X = 4$.}

\subsubsection*{Step 3: the verifier's first challenges}
Once the ciphertexts and commitments are posted, the verifier draws three
coefficient vectors $\beta, \gamma_1, \gamma_2 \in \Z_q^{n}$ and three
scalars $y_1, y_2, y_3$. They must come after the commitments: each of them
compresses many equations into one, and the compression is sound only for
coefficients the voter could not foresee (Steps 4 and 5).

\subsubsection*{Step 4: one linear argument for every linear fact}
Three facts are linear in the committed values: the ciphertexts encrypt
the committed $b$, $\delta = Ms$, and $v = -\delta + 2p - b$ in binary.
Checking each entry would cost $n$ equations per fact. Random coefficients
compress each fact to one equation.
\begin{itemize}
  \item \emph{The link to the ciphertexts.} Everyone computes $C_\beta =
  \prod_k c_k^{\beta_k} = \Enc(\langle \beta, b^{\mathrm{true}} \rangle;\,
  R)$ with $R = \langle \beta, r \rangle$, by the homomorphism. The voter
  proves $C_\beta = (g^{R},\, g^{\langle \beta, b \rangle} h^{R})$ for the
  \emph{committed} $b$. If $b \ne b^{\mathrm{true}}$, the two inner
  products agree for a $1/q$ fraction of the $\beta$ only (Bellare, Garay,
  and Rabin~\cite{bgr98}).
  \item \emph{The anchoring.} The voter proves $\langle \gamma_1,\, \delta -
  Ms \rangle = \sum_k \gamma_{1,k} (\delta_k - (s_j - s_i)) = 0$.
  \item \emph{The binary expansion.} The voter proves $\langle \gamma_2,\,
  \sum_t 2^{t} X^{(t)} + \delta - 2p + b \rangle = 0$.
\end{itemize}
A non-zero vector is orthogonal to a random $\gamma$ with probability
$1/q$, so the two field equations stand for all $2n$ entry-wise ones.

All three facts, together with ``the five commitments open to these
vectors'', are one statement: knowledge of
\[
  w = (s, \delta, b, p, X, \rho_s, \rho_\delta, \rho_b, \rho_p, \rho_X, R)
  \in \Z_q^{W}, \qquad W = m + 3n + Ln + 6,
\]
with $\varphi(w) = T$ for the linear map $\varphi : \Z_q^{W} \to \G^{7}
\times \Z_q^{2}$ whose nine rows are
\begin{align*}
  &\com(s; \rho_s),\ \com(\delta; \rho_\delta),\ \com(b; \rho_b),\ \com(p;
  \rho_p),\ \com(X; \rho_X),\ g^{R},\ g^{\langle \beta, b \rangle} h^{R},\\
  &\langle \gamma_1, \delta - Ms \rangle,\quad \textstyle\langle \gamma_2,
  \sum_t 2^{t} X^{(t)} + \delta - 2p + b \rangle,
\end{align*}
and target $T = (C_s, C_\delta, C_b, C_p, C_X, C_{\beta,1}, C_{\beta,2}, 0,
0)$. Every row is linear in $w$, so the protocol of
Section~\ref{sec:linsigma} proves it: the voter announces $A = \varphi(a)$
for random $a \in \Z_q^{W}$, receives $c$, and answers $z = a + c w$; the
verifier checks $\varphi(z) = A + c T$, which is seven group equations and
two field equations. Because one response vector $z$ serves all nine rows,
the $b$ in the link row is the $b$ that opens $C_b$, and likewise for the
other shared vectors.

\subsubsection*{Step 5: three zero arguments for the products}
The remaining facts are entry-wise products, which the linear argument
cannot express. Each is one Bayer--Groth zero argument
(Section~\ref{sec:zero}) on the commitments of Step 2:
\begin{enumerate}[nosep]
  \item $b \circ \delta - p = 0$: rows $(C_b, C_p)$ against $(C_\delta,
  \com(-\mathbf{1}; 0))$, at $y_1$;
  \item $X \circ (X - \mathbf{1}) = 0$: row $C_X$ against $C_X \cdot
  \com(-\mathbf{1}; 0)$, at $y_2$;
  \item $b \circ (b - \mathbf{1}) = 0$: row $C_b$ against $C_b \cdot
  \com(-\mathbf{1}; 0)$, at $y_3$.
\end{enumerate}
The verifier forms the commitments to $X - \mathbf{1}$ and $b - \mathbf{1}$
itself, by the homomorphism. Each argument proves the single scalar
equation $\sum_j y^{j} (\text{entry } j) = 0$, a polynomial in $y$ whose
coefficients are the entries; for a random $y$ a non-zero vector passes
with probability at most $(\text{length})/q$ (Schwartz--Zippel). Each
argument has its own announcement and its own second challenge $x_i$.

\subsubsection*{Step 6: answer and verify}
After the announcements of Steps 4 and 5, the verifier sends $c$ and $x_1,
x_2, x_3$, and the voter sends the responses. The verifier
(\rocq{verify\_ballot2}) checks that there are $n$ ciphertexts, the linear
argument, and the three zero arguments; if all pass, the $c_k$ are added
to the encrypted margins.

\subsubsection*{Step 7: why accepted means valid}
The checks combine as follows (Theorem~\ref{thm:b2sound}):
\begin{enumerate}[nosep]
  \item the bit arguments make every $b_k$ and every $X_j$ a bit, so $v_k$
  is an integer in $[0, 2^{L})$, and the field is far larger than $2^{L}$,
  so nothing wraps around;
  \item the product argument and the binary row give $v = -\delta + 2 b
  \circ \delta - b$; hence $b_k = 1$ forces $\delta_k = v_k + 1 \ge 1$ and
  $b_k = 0$ forces $\delta_k = -v_k \le 0$;
  \item the anchoring row gives $\delta_{ij} = s_j - s_i$ for one $s$;
  \item together, $b_{ij} = [s_i < s_j]$, and the link row makes the
  ciphertexts encrypt exactly this $b$.
\end{enumerate}

\subsubsection*{Step 8: removing the interaction}
In the non-interactive ballot the ten challenges are hashes
(\rocq{challenges\_of}):
\begin{itemize}[nosep]
  \item $\beta, \gamma_1, \gamma_2, y_1, y_2, y_3$ from the ciphertexts and
  the five commitments;
  \item $c$ from those and the linear announcement $A$;
  \item $x_i$ from those and the announcement of the $i$-th zero argument.
\end{itemize}
Each challenge is hashed only after everything it must bind is fixed. In
particular $\beta$ and $\gamma$ depend on the commitments, so the voter
cannot choose $b$, $\delta$, or $X$ after seeing them. The honest voter
computes the ballot in three passes: the first fixes the ciphertexts and
commitments, hence $\beta, \gamma, y$; the second fixes $A$, which depends
on $\beta$ and $\gamma$, hence $c$; the third computes the responses
(\rocq{nizk\_encrypt\_ballot2}).

\subsubsection*{Step 9: what it costs}
One linear argument and three zero arguments, each linear in $n$, replace
the $2m^3$ Chaum--Pedersen transcripts of the first protocol. At $m = 17$
the batched ballot is $352$\,KB, takes $5.5$\,s to prove and $7.8$\,s to
verify, against $6.2$\,MB, $130$\,s, and $68$\,s for the first.
Completeness, special soundness, and SHVZK are machine-checked, and the
verifier is extracted to OCaml (Section~\ref{sec:mech}).

\subsection{The simplified batched protocol}\label{sec:steps3}
The batched protocol proves four facts about committed vectors with four
arguments and ten challenges. Two of the facts are linear, and a verifier
who holds Pedersen commitments can check a linear fact by multiplying the
commitments instead of asking for a proof; the two remaining facts are
Hadamard products, and several of those can share one zero argument. The
simplified protocol does both. It is the ballot of the paper and of
\rocq{Frontend/SchulzeSimple.v}; the batched protocol above is its first
form, kept here because it shows each fact with its own argument. The
statement is unchanged, so Step~1 of the batched protocol applies
verbatim: $\delta = Ms$, $p = b \circ \delta$, $v = -\delta + 2p - b =
\sum_t 2^t X^{(t)}$, and $b, X^{(0)}, \ldots, X^{(L-1)}$ are bit vectors.
Section~\ref{sec:ballot2sd} gives every number for $A > B > C$ and the
proofs; this walk gives the protocol.

\subsubsection*{Step 1: let the key compute $\delta = Ms$}
A Pedersen commitment is a product of powers of the generators, so
substituting a public linear map into the generators substitutes it into
the committed vector. For the difference matrix $M$, define the
\emph{derived key} $g'_i = \prod_k g_k^{M_{ki}}$, one generator per
candidate, computed by anyone from the key $\mathbf{g}$ and the rule for
$M$. Then
\[
  \com_{\mathbf{g}'}(s; \rho) = h_1^{\rho} \prod_i {g'_i}^{s_i}
  = h_1^{\rho} \prod_k g_k^{\sum_i M_{ki} s_i} = \com_{\mathbf{g}}(Ms; \rho),
\]
so a commitment to $s$ under $\mathbf{g}'$ is, as a group element, a
commitment to $\delta = Ms$ under $\mathbf{g}$ with the same randomness
(\rocq{com\_gd}). The voter posts this one element as $C_\delta$; there
is no $C_s$, and the relation $\delta = Ms$ holds by construction. The
linear argument will still ask the voter to open $C_\delta$ under
$\mathbf{g}'$, so that the extractor recovers $s$.

\ex{For $A, B, C$ with $\mathbf{g} = (34, 8, 27, 14, 36, 3, 17, 42, 12)$ in
the toy group, $g'_A = g_{BA}\, g_{CA}\, g_{AB}^{-1} g_{AC}^{-1} = 32$,
$g'_B = 8$, $g'_C = 9$ (Section~\ref{sec:ballot2sd}, Step~1), and
$\com_{\mathbf{g}'}((0, 1, 2); 14) = 32^{14} \cdot 8 \cdot 9^2 = 27$, the
same element as $\com_{\mathbf{g}}(\delta; 14)$ in
Section~\ref{sec:walk}.}

\subsubsection*{Step 2: commit under one key}
The voter posts the $n$ ciphertexts $c_k = \Enc(b_k; r_k)$ and $L + 3$
commitments, all under the one key $(h_1, \mathbf{g})$ of $n$ generators:
$C_\delta$ from Step~1, $C_b = \com(b; \rho_b)$, $C_p = \com(p;
\rho_p)$, and one commitment $C_{X^{(t)}} = \com(X^{(t)};
\rho_{X^{(t)}})$ per bit level $t < L$. The batched protocol packed the
bits into one vector of length $Ln$ under a second, longer key; keeping
the levels apart costs $L - 1$ extra group elements and puts every
committed vector under the same key, which the next two steps need.

\subsubsection*{Step 3: let the verifier check the binary expansion}
The expansion $v = \sum_t 2^t X^{(t)}$ with $v = -\delta + 2p - b$ says
that the vector $\delta + b - 2p + \sum_t 2^t X^{(t)}$ is zero. Every
term is committed, and the commitment is homomorphic, so the verifier
multiplies:
\[
  C_0 = C_\delta \cdot C_b \cdot C_p^{-2} \cdot \prod_{t < L} C_{X^{(t)}}^{\,2^t}
  = \com\Big(\delta + b - 2p + \textstyle\sum_t 2^t X^{(t)};\ \rho^*\Big),
  \qquad
  \rho^* = \rho_\delta + \rho_b - 2\rho_p + \textstyle\sum_t 2^t \rho_{X^{(t)}} .
\]
For an honest voter the committed vector is zero, so $C_0 = h_1^{\rho^*}$,
and the voter proves knowledge of $\rho^*$: one Schnorr row in the linear
argument. By binding, this proves the relation entry by entry, where the
batched protocol proved it in one random direction $\gamma_2$ and needed
the $n$ one-coordinate variants of $\gamma_2$ to separate the entries.

\ex{With the commitments of Section~\ref{sec:ballot2sd}, $C_0 = 27
\cdot 4 \cdot 17^{-2} \cdot 42 \cdot 21^{2} = 8 = 32^{19}$, and $\rho^*
= 14 + 3 - 40 + 9 + 10 \equiv 19$.}

\subsubsection*{Step 4: the first challenges}
Once the ciphertexts and the $L + 3$ commitments are posted, the verifier
draws $\beta \in \Z_q^{n}$ for the link to the ciphertexts, as before,
and two scalars $y$ and $z$ for the zero argument. Both sides compute
$C_\beta = \prod_k c_k^{\beta_k}$ and $C_0$. There are no $\gamma$
vectors: the two field rows they served are gone.

\subsubsection*{Step 5: the linear argument, five rows}
The witness shrinks to $w = (s, b, \rho_\delta, \rho_b, R, \rho^*)$, of
length $m + n + 4$, and the linear map $\varphi'$ into $\G^5$ has the rows
\[
  \com_{\mathbf{g}'}(s; \rho_\delta),\quad \com(b; \rho_b),\quad g^{R},\quad
  g^{\langle \beta, b \rangle} h^{R},\quad h_1^{\rho^*},
\]
with target $(C_\delta, C_b, C_{\beta,1}, C_{\beta,2}, C_0)$. The first row
recovers $s$, the second $b$, the third and fourth are the link of the
batched protocol, and the fifth is Step~3. The protocol is the one of
Section~\ref{sec:linsigma}: announce $\varphi'(a)$, receive $c$, answer $u
= a + c\,w$. The vectors $\delta$, $p$, and $X^{(t)}$ are no longer in
the witness; they enter only through the zero argument.

\subsubsection*{Step 6: one zero argument for everything quadratic}
The four Hadamard facts, $b \circ \delta = p$, $X^{(t)} \circ (X^{(t)} -
\mathbf{1}) = 0$ for each $t$, and $b \circ (b - \mathbf{1}) = 0$, are
$L + 2$ relations of the kind the zero argument proves, each a sum of
products that must vanish coordinate by coordinate. One argument with $L
+ 3$ rows proves them all, provided no two relations can cancel each
other. The weights $z, z^2, \ldots, z^{2+L}$ prevent that:
\begin{center}
\begin{tabular}{@{}lll@{}}
\toprule
row & left & right \\
\midrule
$1$ & $z\, b$ & $\delta$ \\
$2$ & $z\, p$ & $-\mathbf{1}$ \\
$3 + t$ & $z^{2+t} X^{(t)}$ & $X^{(t)} - \mathbf{1}$ \\
$3 + L$ & $z^{2+L}\, b$ & $b - \mathbf{1}$ \\
\bottomrule
\end{tabular}
\end{center}
The claim $\sum_i \text{left}_i \star_y \text{right}_i = 0$ reads
\[
  z\,\big(b \star_y \delta - p \star_y \mathbf{1}\big)
  + \sum_{t < L} z^{2+t}\, X^{(t)} \star_y (X^{(t)} - \mathbf{1})
  + z^{2+L}\, b \star_y (b - \mathbf{1}) = 0,
\]
a polynomial in $z$ whose coefficients are the $L + 2$ relations at $y$
(rows $1$ and $2$ share the weight $z$ because only their difference
vanishes). If some relation fails, the polynomial is non-zero and has at
most $L + 1$ roots, so a random $z$ catches it except with probability
$(L + 1)/q$; and each relation at $y$ is itself a polynomial in $y$ whose
coefficients are the coordinates, as in the batched protocol. The
verifier builds all $2(L + 3)$ commitments of the table from the posted
ones, $C \mapsto C^{z^i}$ for the weights and $C \mapsto C \cdot
\com(-\mathbf{1}; 0)$ for the subtractions, so the argument costs the
voter one announcement and one response of about $2n$ field elements,
where the three arguments of the batched protocol cost $(2L + 4)n$.

\subsubsection*{Step 7: answer and verify}
The verifier sends $c$ and $x$, the voter answers, and the verifier
(\rocq{verify\_ballot3}) checks that there are $n$ ciphertexts and $L$
bit commitments, the five equations of the linear argument, and the zero
argument on the rows of Step~6. If all pass, the $c_k$ join the
encrypted margins.

\subsubsection*{Step 8: why accepted means valid}
The extractor (\rocq{ballot3\_sound}, Section~\ref{sec:ballot2sd},
Step~7) combines the checks in a different order from the batched
protocol, because two facts now come from binding rather than from a
proof:
\begin{enumerate}[nosep]
  \item the linear argument gives $s$, $b$, and the link, so every
  ciphertext encrypts its committed bit (as in the batched protocol);
  \item one accepted zero argument gives openings $p$ and $X^{(t)}$ of
  $C_p$ and the $C_{X^{(t)}}$; every other accepted $(y, z)$ opens the same
  vectors, by binding, so the weighted claim holds at all of them;
  \item $L + 2$ values of $z$ separate the rows and $n$ values of $y$ the
  coordinates: $p = b \circ \delta$, and $b$ and every $X^{(t)}$ are bit
  vectors;
  \item $C_\delta$ opens to $Ms$ by the derived key, and $C_0 = h_1^{\rho^*}$
  opens to the zero vector, so $\delta = Ms$ and $v = \sum_t 2^t X^{(t)}$
  entry by entry, both by binding;
  \item the decoding of the batched protocol's Step~7 then gives $b_{ij}
  = [s_i < s_j]$.
\end{enumerate}

\subsubsection*{Step 9: removing the interaction}
Five challenges instead of ten (\rocq{challenges\_of3}): $\beta$, $y$,
and $z$ from the ciphertexts and the $L + 3$ commitments; $c$ from those
and the linear announcement; $x$ from those and the zero argument's
announcement. The honest voter still needs three passes, since the linear
announcement depends on $\beta$ and the zero argument's on $y$ and $z$
(\rocq{nizk\_encrypt\_ballot3}).

\subsubsection*{Step 10: what it costs}
The ballot is $2n + 3L + 17$ group elements, of which $2n$ are the
ciphertexts, and $3n + m + 7$ field elements. At $m = 17$ that is
$180$\,KB against $352$\,KB for the batched protocol and $6.2$\,MB for
the first, with $0.7$\,s to prove and $1.6$\,s to verify against
$4.2$\,s and $5.7$\,s (Section~\ref{sec:ballot2s}). The ciphertexts are
now most of the ballot, $148$ of the $180$\,KB. Completeness, special
soundness, SHVZK, and the end-to-end theorem are machine-checked
(\rocq{ballot3\_valid}, \rocq{ballot3\_sound}, \rocq{ballot3\_zk},
\rocq{schulze\_end\_to\_end3}), and the Gentoo elections run through it
take a second or two per ballot (Section~\ref{sec:ballot2s}).

\section{The first ballot: position-anchored membership proofs}\label{sec:ballot1}

\subsection{Encoding a ranking}
A ranking with ties is a \emph{position vector} $s \in \{0, \ldots, m\}^m$:
$s_i$ is the position of candidate $i$, smaller is better, equal positions
are ties, and $s_i = m$ means ``unranked'' (ranked last). Any such vector
is a valid ranking, so no transitivity condition is needed. The ballot's
contribution to the margin $D_{ij}$ is the bit
\[
  b_{ij} = [\, s_i < s_j \,] = [\, \delta_{ij} \ge 1 \,], \qquad
  \delta_{ij} = s_j - s_i \in \{-m, \ldots, m\}.
\]

\ex{Ballot 1, $A > B = C$, is $s = (0, 1, 1)$. Then $\delta_{AB} =
\delta_{AC} = 1$, $\delta_{BA} = \delta_{CA} = -1$, $\delta_{BC} =
\delta_{CB} = 0$; the bits are $b_{AB} = b_{AC} = 1$ and all others $0$. A
tie contributes nothing to either margin. Ballot 2, $B > C > A$, is $s =
(2, 0, 1)$ with bits $b_{BC} = b_{BA} = b_{CA} = 1$.}

\subsection{What the voter publishes}
The voter publishes:
\begin{enumerate}[nosep]
  \item[(P)] for every candidate $i$, $\Enc(s_i)$ with a membership proof
  $s_i \in \{0, \ldots, m\}$ ($m + 1$ branches);
  \item[(B)] for every ordered pair $(i, j)$, $\Enc(b_{ij})$ with a
  membership proof $b_{ij} \in \{0, 1\}$ (2 branches) and a proof of
  consistency with the positions:
  \begin{equation}\label{eq:cons}
    (b_{ij} = 1 \wedge \delta_{ij} \in \{1, \ldots, m\})
    \;\vee\;
    (b_{ij} = 0 \wedge \delta_{ij} \in \{-m, \ldots, 0\}).
  \end{equation}
\end{enumerate}
The verifier computes $\Enc(\delta_{ij}) = \Enc(s_j) / \Enc(s_i)$ by the
homomorphism; the voter never supplies it. This is the \emph{anchoring}. If
$\Enc(\delta_{ij})$ were an input, a voter could post $\delta_{AB} =
\delta_{BA} = 1$ and prove $b_{AB} = b_{BA} = 1$: a ballot that counts $A >
B$ and $B > A$ at once. Anchored to one position vector, $\delta_{BA} =
-\delta_{AB}$ automatically.

\ex{Why a cyclic ballot is impossible. Suppose a voter wants to publish
$b_{AB} = b_{BC} = b_{CA} = 1$, a cycle that no ranking produces and that
would add one to each of $D_{AB}$, $D_{BC}$, $D_{CA}$. The three cells
need $\delta_{AB} \ge 1$, $\delta_{BC} \ge 1$ and $\delta_{CA} \ge 1$,
i.e.\ $s_A < s_B < s_C < s_A$, which no integers satisfy. The voter cannot
use different positions for different cells, because each position is
encrypted once and the verifier computes every $\Enc(\delta_{ij})$ from
the same two position ciphertexts. With $A > B > C$, $s = (0, 1, 2)$, the
cell $(C, A)$ has $\delta_{CA} = -2$: the claim $b_{CA} = 1$ would need
the shifted plaintext $7 \cdot 1 - 2 = 5$ to be in $\{-3, \ldots, 0\}
\cup \{8, 9, 10\}$, which it is not, and only $b_{CA} = 0$, with $t =
-2$, can be proved. Theorem~\ref{thm:b1sound} is the general statement:
every accepted bit is $[s_i < s_j]$ for one vector $s$, and ``$<$'' on
integers has no cycles.}

\paragraph{The shift trick} Condition~\eqref{eq:cons} is two membership
statements joined by ``or'', which would be a disjunction of disjunctions.
It is implemented as \emph{one} membership proof on the ciphertext
\[
  \Enc(b_{ij})^{K} \cdot \Enc(\delta_{ij}), \qquad K = 2m + 1,
\]
whose plaintext is $t = K b_{ij} + \delta_{ij}$, with message set $\{-m,
\ldots, 0\} \cup \{K + 1, \ldots, K + m\}$. Since $b \in \{0, 1\}$ (by its
own proof) and $\delta \in [-m, m]$ (by the position proofs), $t \in \{-m,
\ldots, 0\}$ exactly when $b = 0$ and $\delta \le 0$, and $t \in \{K+1,
\ldots, K+m\}$ exactly when $b = 1$ and $\delta \ge 1$: the two ranges are
separated by the gap $\{1, \ldots, K\}$, which no honest or dishonest $(b,
\delta)$ can straddle in the wrong direction.

\ex{$m = 3$, $K = 7$, message set $\{-3, -2, -1, 0\} \cup \{8, 9, 10\}$.
Honest cells: $(b, \delta) = (1, 1)$ gives $t = 8$; $(0, -1)$ gives $t =
-1$; $(0, 0)$ gives $t = 0$. Dishonest cells: $(1, 0)$ gives $7$, $(1, -1)$
gives $6$, $(0, 1)$ gives $1$, $(0, 2)$ gives $2$: none is in the set, and
the membership proof cannot be made. Table~\ref{tab:ballot1} lists ballot
1.}

\begin{table}[t]
  \centering
  \caption{Ballot 1 in the first protocol: $s = (0,1,1)$, $m = 3$, $K =
  7$. Every row is one ciphertext $\Enc(b_{ij})$ with its 0/1 proof and
  one membership proof on $\Enc(b_{ij})^7 \Enc(\delta_{ij})$. The last two
  rows are cheating attempts.}
  \label{tab:ballot1}
  \begin{tabular}{@{}lcccl@{}}
    \toprule
    pair & $b_{ij}$ & $\delta_{ij}$ & $t = 7 b_{ij} + \delta_{ij}$ & membership \\
    \midrule
    $(A,B)$ & 1 & $1$ & $8$ & $\in \{8,9,10\}$ \\
    $(A,C)$ & 1 & $1$ & $8$ & $\in \{8,9,10\}$ \\
    $(B,A)$ & 0 & $-1$ & $-1$ & $\in \{-3,\ldots,0\}$ \\
    $(C,A)$ & 0 & $-1$ & $-1$ & $\in \{-3,\ldots,0\}$ \\
    $(B,C)$ & 0 & $0$ & $0$ & $\in \{-3,\ldots,0\}$ \\
    $(C,B)$ & 0 & $0$ & $0$ & $\in \{-3,\ldots,0\}$ \\
    \midrule
    $(B,A)$ & 1 & $-1$ & $6$ & rejected \\
    $(B,C)$ & 1 & $0$ & $7$ & rejected \\
    \bottomrule
  \end{tabular}
\end{table}

A ballot thus has $m + m^2$ ciphertexts (the implementation keeps the $m$
diagonal cells, which encrypt $0$), $m$ membership proofs with $m+1$
branches and $m^2$ cells with a 2-branch and a $(2m+1)$-branch proof: about
$2m^3$ Chaum--Pedersen transcripts. All proofs of a ballot share one
Fiat--Shamir challenge, the hash of the voter's credential and of all
ciphertexts and announcements (\rocq{nizk\_encrypt\_ballot},
\rocq{nizk\_verify\_ballot}).

\subsection{Soundness}
\begin{theorem}[\rocq{ballot\_sound}]\label{thm:b1sound}
Let $h = g^x$, let $g$ have order $q$ (\rocq{exp\_injective}), and let the
field have characteristic above $4m + 2$ (\rocq{natF\_injective}). From two
accepting transcripts of a ballot with the same announcements and different
challenges, there is a position vector $s \in \{0, \ldots, m\}^m$ such that
every pair ciphertext decrypts to $g^{[s_i < s_j]}$.
\end{theorem}
\begin{proof}
Every proof of the ballot is an OR-composition of Chaum--Pedersen proofs
with the shared challenge, so the two transcripts give two transcripts of
every proof with distinct challenges, and the extractor of
Section~\ref{sec:sigma} applies to each.

(1) From the position proofs: for every $i$ there are $s_i \in \{0, \ldots,
m\}$ and $r_i$ with $\Enc(s_i; r_i)$ the posted ciphertext. Define $s$ from
these; it is a position vector.

(2) From the 0/1 proofs: $\Enc(b_{ij})$ encrypts some $b_{ij} \in \{0, 1\}$.

(3) From the consistency proofs: the ciphertext $\Enc(b_{ij})^K \cdot
\Enc(s_j) / \Enc(s_i)$, which by the homomorphism encrypts $t = K b_{ij} +
s_j - s_i$ computed in $\Z_q$, encrypts some $t' \in \{-m, \ldots, 0\} \cup
\{K+1, \ldots, K+m\}$. Two plaintexts of one ciphertext under one key are
equal as group elements, $g^{t} = g^{t'}$, hence $t = t'$ in $\Z_q$ because
$g$ has order $q$. Both $t$ and $t'$ are images of integers in $[-m, 3m+1]$,
an interval of length $4m + 1$, so $t = t'$ as integers by the
characteristic bound.

(4) Case $b_{ij} = 0$: $t = \delta_{ij} \in [-m, m]$, and $t \in \{K+1,
\ldots\}$ is impossible since $K + 1 = 2m + 2 > m$; so $\delta_{ij} \le 0$,
i.e.\ $s_i \ge s_j$ and $[s_i < s_j] = 0 = b_{ij}$. Case $b_{ij} = 1$: $t =
K + \delta_{ij} \ge K - m = m + 1 > 0$, so $t \in \{K+1, \ldots, K+m\}$,
hence $\delta_{ij} \ge 1$ and $[s_i < s_j] = 1 = b_{ij}$. In both cases
$\Dec(\Enc(b_{ij})) = g^{b_{ij}} = g^{[s_i < s_j]}$.
\end{proof}
The two side conditions are facts about the Schnorr group, proved once for
every instance in \rocq{Utility/SchnorrOrder.v}; at the instance the
theorems have no hypothesis beyond the transcripts
(\rocq{ballot\_sound\_ins}).

\subsection{Zero-knowledge}
\begin{theorem}[\rocq{ballot\_zk}]
For positions in $[0, m]$, the distribution of the ballot's proofs over the
voter's randomness is a permutation of the distribution output by a
simulator that has the ciphertexts and the challenge only.
\end{theorem}
\begin{proof}
Each proof is simulated by the OR-simulator (all branches simulated), and
the ballot's proofs are independent given the ciphertexts; the development
lifts the per-proof permutations through the product of the $m + 2m^2$
distributions (\rocq{dist\_pair\_perm}, \rocq{dist\_vec\_perm}). The
simulator derives $\Enc(\delta_{ij})$ from the position ciphertexts as the
verifier does, so it needs nothing the verifier does not have.
\end{proof}
The ciphertexts themselves are not simulated: what they reveal is governed
by IND-CPA, an on-paper reduction.

\subsection{Cost}
About $2m^3$ Chaum--Pedersen transcripts, each four group elements and two
field elements in the ballot and a few exponentiations for prover and
verifier. At $m = 17$ (the largest Gentoo election): $6.2$\,MB, $130$\,s to
prove and $68$\,s to verify at the $2048$-bit parameters. This is the cost
the batched ballot removes.

\section{Aggregation}\label{sec:agg}
The encrypted margin matrix is the product of the accepted ballots' pair
ciphertexts, $\Enc(D_{ij}) = \prod_{\text{ballots}} \Enc(b_{ij})$, which by
the homomorphism encrypts the number of accepted voters with $b_{ij} = 1$,
i.e.\ $D_{ij}$. Anyone computes it from the board with no key
(\rocq{aggregate\_margins}, \rocq{cdec\_aggregate}: ``decrypting the
aggregate gives $g^{D_{ij}}$ when every ballot's bit decrypts to
$g^{b_{ij}}$''). Rejected ballots are posted as rejected and left out.

\ex{For the pair $(A, B)$: ballots $1, 4, 5$ have $b_{AB} = 1$ and ballots
$2, 3$ have $0$; the product of the five ciphertexts encrypts $3 = D_{AB}$.
In the toy group, if the five ciphertexts have randomness $r_1, \ldots,
r_5$, the aggregate is $\Enc(3; r_1 + \cdots + r_5)$, and the trustees'
decryption gives $g^3 = 8$.}

\section{The batched ballot: four vector arguments}\label{sec:ballot2}

\subsection{The idea}
The first ballot pays $2m^3$ transcripts because a $\Sigma$-protocol speaks
about one ciphertext at a time, and ``$\delta_{ij}$ is in a range'' is a
disjunction over the range. The batched ballot commits to the whole vectors
$s, \delta, b$ and proves the same validity as three statements about
vectors: a linear relation, an entry-wise product, and ``these entries are
bits''. Each is one argument of Section~\ref{sec:tools} over the whole
vector, with cost linear in the vector length. Positions are
\emph{committed}, not encrypted; only the $m^2$ bits are encrypted, since
only they are aggregated.

\subsection{The statement}
Write $n = m^2$ and index pairs $(i, j)$ by $k = i m + j \in \{0, \ldots, n
- 1\}$. Let $L = \lceil \log_2(m + 1) \rceil$. For a position vector $s$
define
\[
  \delta = M s \quad (\delta_{ij} = s_j - s_i), \qquad
  p = b \circ \delta \quad (p_k = b_k \delta_k), \qquad
  v = -\delta + 2p - b .
\]
Here $M$ is the $n \times m$ \emph{difference matrix}. Each difference
is a fixed linear combination of the positions: $\delta_{ij} = (+1)\cdot s_j
+ (-1)\cdot s_i$, with coefficient $0$ on every other position. A row of
$M$ holds the coefficients of one such combination, so each entry of $M$ is
a single coefficient, $-1$, $0$, or $+1$, and the whole matrix stacks all
$n$ rows so that $\delta = M s$ computes every difference at once. For
three candidates, row $AB$ is $(-1, +1, 0)$, and applied to $s = (s_A,
s_B, s_C)$ it gives
\[
  \delta_{AB} = (-1)\cdot s_A + (+1)\cdot s_B + 0 \cdot s_C = s_B - s_A .
\]
In general row $(i, j)$ has $+1$ in column $j$, $-1$ in column $i$, and
$0$ elsewhere, and the diagonal rows $(i, i)$ are zero. The matrix $M$ is public and does not depend on any ballot: each row
only says which two positions to subtract, not what they are. Anyone who
knows the number of candidates can compute it, and it is the same for
every voter. The verifier builds it itself: the candidate list, and so
$m$, is part of the published election setup, and the rule for $M$ is part
of the protocol that every verifier implements, like its verification
equations. $M$ never appears on the board, and the voter cannot choose it.
What is secret is the voter's $s$, and through it $\delta = Ms$. The
development never builds $M$ as a matrix; the linear map $\varphi$
computes each row as the subtraction $s_j - s_i$.

Computing $M$ for any candidate list takes three steps. First, number the
candidates $0, 1, \ldots, m - 1$ in the order of the published list; this
order also fixes the order of the voter's ciphertexts. Second, for every
ordered pair $(i, j)$, write the row vector of length $m$ that has $-1$ in
entry $i$, $+1$ in entry $j$, and $0$ everywhere else; with $e_i$ the unit
row vector that has $1$ in entry $i$, this row is $e_j - e_i$. When $i =
j$ the $-1$ and the $+1$ fall in the same entry and cancel, so the
diagonal rows are zero without a special case. Third, stack the $m^2$ rows
in the order of the cells, with row $(i, j)$ in position $k = i m + j$. The
result is $M$, and row $k$ applied to $s$ gives $\delta_{ij} = s_j - s_i$.
For the list $A, B, C$, numbered $0, 1, 2$, the pair $(B, C) = (1, 2)$ is
row $k = 1 \cdot 3 + 2 = 5$, the vector $e_2 - e_1 = (0, -1, +1)$. Applied
to $m = 3$, the three steps give the matrix below, written as $M^{\top}$
with rows for $A, B, C$ and columns for the cells $AA, AB, AC, BA, BB, BC,
CA, CB, CC$:
\[
  M^{\top} =
  \begin{pmatrix}
    0 & -1 & -1 & 1 & 0 & 0 & 1 & 0 & 0 \\
    0 & 1 & 0 & -1 & 0 & -1 & 0 & 1 & 0 \\
    0 & 0 & 1 & 0 & 0 & 1 & -1 & -1 & 0
  \end{pmatrix}.
\]
Two voters in the same election share this $M$ and differ only in $s$,
hence in $\delta$. The ranking $A > B > C$ has $s = (0, 1, 2)$ and $C > B
> A$ has $s = (2, 1, 0)$:
\begin{align*}
  A > B > C:&\quad \delta = M\,(0, 1, 2) = (0, 1, 2,\; -1, 0, 1,\; -2, -1, 0),\\
  C > B > A:&\quad \delta = M\,(2, 1, 0) = (0, -1, -2,\; 1, 0, -1,\; 2, 1, 0).
\end{align*}
The reversed ranking negates every difference, as it should; the matrix is
the same in both cases.
If $b_k = 0$ then $v_k = -\delta_k$; if $b_k = 1$ then $v_k = \delta_k - 1$.
So $b_k = [\delta_k \ge 1]$ is equivalent to ``$b_k$ is a bit and $v_k$ is a
non-negative integer'', and since $|\delta_k| \le m$, a non-negative $v_k$
is below $2^L$. The validity statement is therefore
\begin{equation}\label{eq:validity}
\begin{gathered}
  b \in \{0,1\}^n,\qquad
  \delta = M s,\qquad
  p = b \circ \delta,\\
  v = \sum_{t < L} 2^t X^{(t)} \text{ for bit vectors } X^{(0)}, \ldots, X^{(L-1)} \in \{0,1\}^n .
\end{gathered}
\end{equation}
No upper bound on $s$ is needed: the differences are small, so the
positions relative to one candidate are small integers.

The vector $p$ is not needed to define $v$: mathematically $v = -\delta +
2(b \circ \delta) - b$. It is needed by the proofs. The linear argument
(a) below proves only relations that are linear in the committed vectors,
with public coefficients, and $b \circ \delta$ multiplies two secret
vectors. So the ballot commits to $p$ as a vector of its own, which makes
the definition of $v$ linear in $(\delta, p, b)$, and proves the one
quadratic fact $p = b \circ \delta$ with a zero argument. Bit-ness,
$b_k(b_k - 1) = 0$ and $X_k(X_k - 1) = 0$, is quadratic for the same
reason and also goes to zero arguments; everything else ($\delta = Ms$, the
binary expansion of $v$, and the link to the ciphertexts) is linear and goes
to the single linear argument. This is the usual treatment of
multiplication in proofs for arithmetic circuits: every product gets a
committed output of its own.

\ex{Ballot 1, $s = (0, 1, 1)$, $m = 3$, $n = 9$, $L = 2$, cells in the
order $AA$, $AB$, $AC$, $BA$, $BB$, $BC$, $CA$, $CB$, $CC$:
\begin{align*}
  \delta &= (0, 1, 1,\; -1, 0, 0,\; -1, 0, 0), \\
  b &= (0, 1, 1,\; 0, 0, 0,\; 0, 0, 0), \\
  p &= (0, 1, 1,\; 0, 0, 0,\; 0, 0, 0), \\
  v &= (0, 0, 0,\; 1, 0, 0,\; 1, 0, 0), \\
  X^{(0)} &= v, \qquad X^{(1)} = 0 .
\end{align*}
Check cell $AB$: $v = -1 + 2 \cdot 1 - 1 = 0$; cell $BA$: $v = 1 + 0 - 0 =
1$. Now the cheating attempts of Table~\ref{tab:ballot1}: $b_{BA} = 1$ with
$\delta_{BA} = -1$ forces $p_{BA} = -1$ and $v_{BA} = 1 - 2 - 1 = -2$;
$b_{BC} = 1$ with $\delta_{BC} = 0$ forces $v_{BC} = -1$. In $\Z_{23}$
these are $21$ and $22$, neither of which is $X^{(0)} + 2 X^{(1)}$ for bits
$X$. So no witness satisfies~\eqref{eq:validity}, and the cheating voter
fails at the bit argument instead of at a membership proof.}

\subsection{Where the statement comes from}
The vectors $\delta$, $p$, and $v$ are what is left after the validity
condition $b_{ij} = [s_i < s_j]$ is rewritten in the only three shapes the
tools of Section~\ref{sec:tools} can check about committed vectors: linear
relations (the linear argument), entry-wise products, and ``every entry is
$0$ or $1$'' (both by zero arguments). Each step of the rewriting is forced
by one limitation of these tools.

The first step anchors the differences. As in the first ballot, every
pair must use differences of one position vector, or a voter could choose
$\delta_{AB} = \delta_{BA} = 1$. The differences are a linear function of
$s$, so the anchoring is the linear relation $\delta = Ms$.

The second step turns the comparison into a single non-negativity claim.
The condition $b = [\delta \ge 1]$ is a case split: if $b = 1$ then
$\delta - 1 \ge 0$, and if $b = 0$ then $-\delta \ge 0$. Both cases say
that some quantity is non-negative, and for a bit $b$ the standard way to
write ``if $b$ then $x$ else $y$'' is $b x + (1 - b) y$. With $x = \delta -
1$ and $y = -\delta$,
\[
  v = b(\delta - 1) + (1 - b)(-\delta) = -\delta + 2 b \delta - b ,
\]
so the condition becomes: $b$ is a bit and $v$ is non-negative.

The third step expresses non-negativity in a field, where elements wrap
around modulo $q$ and ``$\ge 0$'' has no meaning. The voter instead shows
that $v$ is a small integer by writing it in binary, $v = \sum_t 2^t
X^{(t)}$ with bit vectors $X^{(t)}$. Honest values are bounded, $v = -\delta
\le m$ when $b = 0$ and $v = \delta - 1 \le m - 1$ when $b = 1$, so $L =
\lceil \log_2(m + 1) \rceil$ bits suffice. The expansion is linear in $X$,
and bit-ness is the third shape.

The fourth step names the product. The term $b\delta$ multiplies two secret
vectors, which the linear argument cannot express, so the voter commits to
$p = b \circ \delta$ separately, as explained above. This makes $v =
-\delta + 2p - b$ linear, and the product argument proves $p = b \circ
\delta$.

Why this is enough: suppose every check passes, so that, after rewinding,
the committed vectors satisfy all the relations and the ciphertexts
encrypt the committed $b$. Take any pair $k = (i, j)$. If $b_k = 0$ then
$v_k = -\delta_k$ is in $[0, 2^L)$, so $\delta_k \le 0$, i.e.\ $s_i \ge
s_j$, and the bit is the correct $[s_i < s_j] = 0$. If $b_k = 1$ then $v_k
= \delta_k - 1$ is in $[0, 2^L)$, so $\delta_k \ge 1$, i.e.\ $s_i < s_j$,
and the bit is the correct $1$. Since $b$ is a bit there is no third case,
so every encrypted bit is the comparison bit of one position vector, which
is exactly a valid ballot.

Every check is needed. Table~\ref{tab:necessary} shows, for each one, an
invalid ballot that passes without it. The examples are in $\Z_{23}$;
the second and third rows start from the ballot $A > B > C$, where
$\delta_{BA} = -1$.

\begin{table}[ht]
  \centering
  \caption{Removing any one check admits an invalid ballot.}
  \label{tab:necessary}
  \begin{tabularx}{\textwidth}{@{}l>{\raggedright\arraybackslash}X@{}}
    \toprule
    removed check & an invalid ballot that then passes \\
    \midrule
    $\delta = Ms$ & commit to $\delta_{AB} = \delta_{BA} = 1$ and set
    $b_{AB} = b_{BA} = 1$: the ballot counts $A > B$ and $B > A$ \\
    $p = b \circ \delta$ & claim $b_{BA} = 1$ and commit to $p_{BA} = 0$
    instead of $-1$: then $v_{BA} = 1 + 0 - 1 = 0$, which looks
    non-negative \\
    $X$ is a bit vector & claim $b_{BA} = 1$ with the honest $p_{BA} =
    -1$: then $v_{BA} = 1 - 2 - 1 = -2 \equiv 21$, and $X^{(0)}_{BA} = 21$
    expands it \\
    $b$ is a bit vector & on the tied pair $B, C$ of the ballot $A > B = C$,
    where $\delta_{BC} = 0$, take $b_{BC} = -1$: then $p_{BC} = 0$ and
    $v_{BC} = 1$, which passes the range check, and the ballot subtracts
    one from the margin $D_{BC}$ \\
    the ciphertext link & commit to the honest $b$ and encrypt different
    bits \\
    \bottomrule
  \end{tabularx}
\end{table}

Two subtleties are handled by the formal proof. First, no range check on
$s$ is needed. From $v$ being small, every difference $s_j - s_i$ is a
small integer, so the positions relative to the first candidate,
$\pos_i = 2^L + (s_i - s_0)$, are natural numbers, and any vector of
natural numbers is a valid ranking with ties. Second, the field must be
large enough. ``$\delta_k \ge 1$'' is a statement about integers, while
the checks hold modulo $q$, and in a small field a negative value could
wrap around and look small. Theorem~\ref{thm:b2sound} therefore assumes a
characteristic above $3 \cdot 2^L + 2$, so that every value involved stays
below $q$; the bound came out of the formal proof. The case analysis above
is the readable version of \rocq{ballot2\_sound}, which checks every step
in Rocq, including both subtleties.

\subsection{The ballot}
Public: the election key $(g, h)$; a commitment key $h_1, \mathbf{g}_1 \in
\G^{n}$ and a longer one $h_1, \mathbf{g}_2 \in \G^{Ln}$. The voter posts:
\begin{itemize}[nosep]
  \item the ciphertexts $c_k = \Enc(b_k; r_k)$, $k < n$;
  \item commitments $C_s = \com(s; \rho_s)$, $C_\delta$, $C_b$, $C_p$ over
  $\mathbf{g}_1$ (with $s$ padded to length $n$), and $C_X = \com(X;
  \rho_X)$ over $\mathbf{g}_2$, where $X$ is the concatenation of
  $X^{(0)}, \ldots, X^{(L-1)}$;
  \item three arguments, (a), (b), (c) below.
\end{itemize}

\paragraph{(a) The linear argument} The verifier draws coefficient vectors
$\beta, \gamma_1, \gamma_2 \in \Z_q^n$ \emph{after} the ciphertexts and
commitments are posted. Everyone computes the ciphertext combination
\[
  C_\beta = \prod_{k} c_k^{\beta_k} = \Enc\big(\langle \beta, b^{\mathrm{true}}\rangle;\; \langle \beta, r \rangle\big),
\]
where $b^{\mathrm{true}}$ is the vector of plaintexts. The voter proves
knowledge of
\[
  w = (s, \delta, b, p, X, \rho_s, \rho_\delta, \rho_b, \rho_p, \rho_X, R)
  \in \Z_q^{W}, \qquad W = m + 3n + Ln + 6,
\]
with $\varphi_{\beta, \gamma}(w) = T$, where $\varphi$ is the linear map
into $\G^7 \times \Z_q^2$ with rows
\begin{align*}
  &\com(s; \rho_s),\; \com(\delta; \rho_\delta),\; \com(b; \rho_b),\; \com(p; \rho_p),\; \com(X; \rho_X),\;
  g^{R},\; g^{\langle \beta, b \rangle} h^{R}, \\
  &\langle \gamma_1,\; \delta - M s \rangle,\qquad
  \big\langle \gamma_2,\; \textstyle\sum_t 2^t X^{(t)} + \delta - 2p + b \big\rangle,
\end{align*}
and $T = (C_s, C_\delta, C_b, C_p, C_X, C_{\beta,1}, C_{\beta,2}, 0, 0)$.
Every row is linear in $w$ (the commitments are exponentials of linear
forms, the inner products are linear forms), so Section~\ref{sec:linsigma}
applies. The rows say: the five commitments open to $(s, \delta, b, p, X)$;
$C_\beta$ encrypts $\langle \beta, b \rangle$ with randomness $R$ for the
\emph{committed} $b$; and the two field identities hold in the direction
$\gamma_1$, $\gamma_2$.

\paragraph{(b) The product argument} A zero argument with rows $(b, p)$
against $(\delta, -\mathbf{1})$ at challenge $y_1$: coordinate relations
$b_k \delta_k - p_k = 0$.

\paragraph{(c) The two bit arguments} A zero argument with the single row
$X$ against $X - \mathbf{1}$ at challenge $y_2$, and one with $b$ against
$b - \mathbf{1}$ at $y_3$; the verifier derives the commitments to $X -
\mathbf{1}$ and $b - \mathbf{1}$ as $C_X \cdot \com(-\mathbf{1}; 0)$ and
$C_b \cdot \com(-\mathbf{1}; 0)$.

\paragraph{Verification} (\rocq{verify\_ballot2}) Check that there are $n$
ciphertexts; check the linear argument, $\varphi_{\beta,\gamma}(z) = A + c
\cdot T$, which unfolds to seven group equations and two field equations;
check the three zero arguments. All challenges are drawn by the verifier
in the interactive version and by hashing in the non-interactive one.

\paragraph{Fiat--Shamir} (\rocq{challenges\_of}, \rocq{nizk\_encrypt\_ballot2})
The challenges are derived in this order, each by a tagged query to the
oracle: $\beta, \gamma_1, \gamma_2$ (one query per coordinate) and $y_1,
y_2, y_3$ from the ciphertexts and the commitments; $c$ from those and the
announcement $A$ of the linear argument; $x_i$ from those and the
announcement of the $i$-th zero argument. The order matters: a voter who
could choose $\beta$ could answer the link with a committed $b$ that is not
the plaintext vector. The honest voter computes the ballot in three
passes: the first fixes the ciphertexts and commitments and hence $\beta,
\gamma, y$ (and the zero arguments' announcements, which do not depend on
$x$); the second fixes $A$, which depends on $\beta, \gamma$, and hence
$c$; the third answers. Ten oracle queries in all; the honest ballot is
accepted (\rocq{nizk\_ballot2\_valid}).

\subsection{Why the random linear combination is sound}
The link is the small-exponent batch test of Bellare, Garay, and
Rabin~\cite{bgr98}, which shuffle arguments use to tie ciphertexts to
committed vectors~\cite{tereliuswikstrom,bayergroth}.
Suppose the committed $b$ differs from the plaintext vector
$b^{\mathrm{true}}$ in coordinate $k$. The link says $C_\beta$ encrypts
$\langle \beta, b \rangle$; but $C_\beta$ encrypts $\langle \beta,
b^{\mathrm{true}} \rangle$. The two agree iff $\beta_k (b_k -
b^{\mathrm{true}}_k) = \sum_{j \ne k} \beta_j (b^{\mathrm{true}}_j - b_j)$,
which for fixed other coordinates holds for exactly one value of
$\beta_k$. So a cheating voter answers the link for a $1/q$ fraction of
the coefficient vectors. The mechanised statement avoids probabilities by
asking for the link at a base $\beta$ and at a $\beta'$ differing only in
coordinate $k$: dividing the two links,
\[
  c_k^{\beta_k - \beta'_k} = \Enc\big(b_k (\beta_k - \beta'_k);\; R - R'\big),
\]
and since $\G$ is a vector space over $\Z_q$, raising both sides to the
power $(\beta_k - \beta'_k)^{-1}$ gives $c_k = \Enc(b_k; (R - R')/(\beta_k
- \beta'_k))$, so $\Dec(c_k) = g^{b_k}$ (\rocq{beta\_bits},
\rocq{root\_unique}). Note that this needs no hypothesis on the order of
$g$, unlike step (3) of Theorem~\ref{thm:b1sound}: nothing is compared in
the exponent. The same two-vector trick on $\gamma_1$ and $\gamma_2$ turns
the two field identities, which hold only in the direction $\gamma$, into
entry-wise identities (\rocq{gamma\_rows}): if $\langle \gamma, e \rangle
= 0$ and $\langle \gamma', e \rangle = 0$ with $\gamma, \gamma'$ differing
only at $k$, then $(\gamma_k - \gamma'_k) e_k = 0$, so $e_k = 0$.

\ex{Toy link. Suppose $n = 2$, the plaintexts are $b^{\mathrm{true}} = (1,
0)$ and a cheating voter commits to $b = (1, 1)$. For $\beta = (3, 5)$,
$C_\beta$ encrypts $3$, but the voter must prove it encrypts $\langle
\beta, b \rangle = 8$: impossible unless $3 = 8$ in $\Z_{23}$. The only
$\beta_2$ for which the voter could answer with $\beta_1 = 3$ is $\beta_2 =
0$, one value in $23$.}

\subsection{Soundness}
\begin{theorem}[\rocq{ballot2\_sound}]\label{thm:b2sound}
Let $h = g^x$, $L \ge 1$, and let the field have characteristic above $3
\cdot 2^L + 2$. Suppose a batched ballot is accepted
\begin{itemize}[nosep]
  \item by the linear argument, on a base tuple $(\beta, \gamma_1,
  \gamma_2)$ with two challenges $c \ne c'$ on the same announcement, and
  for every coordinate $k$ on a tuple differing from the base only in
  $\beta_k$ (likewise only in $\gamma_{1,k}$, likewise only in
  $\gamma_{2,k}$), each with two challenges;
  \item by the product argument on $n$ distinct non-zero values of $y_1$,
  with $2 \cdot 2 + 1 = 5$ challenges each on one announcement;
  \item by the two bit arguments on $Ln$ (resp.\ $n$) distinct non-zero
  values of $y$, with $3$ challenges each.
\end{itemize}
Then either the binding of $(h_1, \mathbf{g}_1)$ or of $(h_1,
\mathbf{g}_2)$ is broken, or there is $\pos \in \{0, \ldots, 2^{L+1}\}^m$
such that every $c_{ij}$ decrypts to $g^{[\pos_i < \pos_j]}$.
\end{theorem}
\begin{proof}
\emph{Step 1, the base witness.} The two challenges at the base tuple give,
by Lemma of Section~\ref{sec:linsigma}, a witness $w$ with
$\varphi_{\beta,\gamma}(w) = T$ (\rocq{lin\_open}). Reading the rows
(\rocq{phi\_target\_open}): $C_s = \com(s; \rho_s)$ and likewise for
$\delta, b, p, X$; $C_\beta = \Enc(\langle \beta, b \rangle; R)$; $\langle
\gamma_1, \delta - Ms \rangle = 0$; $\langle \gamma_2, \sum_t 2^t X^{(t)} +
\delta - 2p + b \rangle = 0$.

\emph{Step 2, the other witnesses open the same vectors.} Each tuple that
varies one coordinate gives, in the same way, a witness $w'$ whose rows
open the same five commitments. If any of the five openings differs from
$w$'s, we have two openings of a Pedersen commitment, a binding break; we
stop with that conclusion (\rocq{open\_agree}). Otherwise every witness
has the same $(s, \delta, b, p, X)$, and we have the link and the two
identities at every tuple.

\emph{Step 3, the committed $b$ is the plaintext vector.} For every $k$,
the links at $\beta$ and at the tuple varying $\beta_k$ give $\Dec(c_k) =
g^{b_k}$ by the division-and-root argument above (\rocq{beta\_bits}).

\emph{Step 4, the linear relations hold entry-wise.} For every $k$, the
identities at $\gamma_1$ and at the tuple varying $\gamma_{1,k}$ give
$\delta_k = s_j - s_i$ for $k = im + j$; likewise $\gamma_2$ gives $v_k :=
\sum_t 2^t X^{(t)}_k = -\delta_k + 2 p_k - b_k$ (\rocq{gamma\_rows}).

\emph{Step 5, the zero arguments.} The $5$ transcripts of the product
argument at each $y_1$ extract openings of $C_b, C_p, C_\delta$ and of
$\com(-\mathbf{1}; 0)$ satisfying $b' \star_{y_1} \delta' + p' \star_{y_1}
u' = 0$; the opening $u'$ of $\com(-\mathbf{1}; 0)$ is $-\mathbf{1}$ and the
others are $b, \delta, p$ unless binding is broken. So $\sum_k y_1^{k}
(b_k \delta_k - p_k) = 0$ at $n$ distinct non-zero $y_1$, hence $p_k = b_k
\delta_k$ for all $k$ (\rocq{zero\_prod\_open}; the Vandermonde argument
is \rocq{PolyRoots}). The bit arguments likewise give $X_j (X_j - 1) = 0$
for all $j < Ln$ and $b_k (b_k - 1) = 0$ for all $k$
(\rocq{zero\_bits\_open}).

\emph{Step 6, decoding the positions} (\rocq{decode\_positions}). Now
every $X_j$ and $b_k$ is $0$ or $1$ in the field, so $v_k = \sum_t 2^t
X^{(t)}_k$ is the image of a natural number $\nu_k < 2^L$. From $p_k = b_k
\delta_k$: if $b_k = 0$ then $\delta_k = -v_k$, and if $b_k = 1$ then
$\delta_k = v_k + 1$. Define, with candidate $0$ as reference, $\pos_i = 2^L
+ (\text{the integer } \delta_{0i})$, i.e.\ $\pos_i = 2^L + \nu_{0i} + 1$
if $b_{0i} = 1$ and $\pos_i = 2^L - \nu_{0i}$ if $b_{0i} = 0$; this is a
natural number in $[0, 2^{L+1}]$ (\rocq{posn}, \rocq{posn\_le}), and in
the field $\pos_i = 2^L + s_i - s_0$. Then $\delta_{ij} = s_j - s_i =
\pos_j - \pos_i$ in the field, and the same case analysis on $b_{ij}$ shows
$b_{ij} = 1 \Leftrightarrow \pos_j - \pos_i = \nu_{ij} + 1 \ge 1$ and
$b_{ij} = 0 \Leftrightarrow \pos_j - \pos_i = -\nu_{ij} \le 0$, where the
integer inequalities follow from the field equalities because every
quantity involved is the image of an integer of absolute value at most $3
\cdot 2^L + 2$ and the characteristic exceeds that. So $b_{ij} = [\pos_i <
\pos_j]$, and by Step 3, $\Dec(c_{ij}) = g^{[\pos_i < \pos_j]}$.
\end{proof}
The theorem is in the \emph{perfect-adversary} form: it assumes the prover
answers every challenge the extractor needs. The usual forking-lemma
argument turns this into a bound on a prover accepted with probability
$\varepsilon$; it is on paper, as for the library's shuffle theorems.

\subsection{Zero-knowledge}
\begin{theorem}[\rocq{ballot2\_zk}]
For positions in $[0, m]$ and $m < 2^L$, and for commitment keys of the
form $g_{1,j} = h_1^{\gamma_{1,j}}$, $g_{2,j} = h_1^{\gamma_{2,j}}$, the
distribution of the honest batched ballot over the voter's randomness (the
$n$ encryption randomnesses, the five commitment randomnesses, the
announcement randomness $a$ of the linear argument and the randomness of
the three zero arguments) is a permutation of the output of a simulator
that has the ciphertexts and the challenges only.
\end{theorem}
\begin{proof}
The simulator draws five exponents $u_s, \ldots, u_X$ uniformly, sets the
commitments to $h_1^{u}$, computes $T$ from the ciphertexts and these
commitments, and runs the simulators of the linear argument
(Section~\ref{sec:linsigma}) and of the three zero arguments on them. The
proof composes three bijections of the randomness
(\rocq{real\_proofs\_perm}, \rocq{com\_vals\_perm}):
\begin{enumerate}[nosep]
  \item For fixed ciphertexts and fixed commitments, the real proofs are
  distributed as the simulated ones: this is \rocq{lin\_shvzk} and
  \rocq{zero\_shvzk} applied to the four arguments one after the other,
  the statements of the zero arguments being rewritten with the
  homomorphism ($\com(X - \mathbf{1}; \rho_X) = C_X \cdot \com(-\mathbf{1};
  0)$).
  \item The real commitments, over the five commitment randomnesses, are
  distributed as five uniform elements of $\langle h_1 \rangle$: with
  $\gamma$ the discrete logarithms of the key, $\com(v; \rho) = h_1^{\rho +
  \langle \gamma, v \rangle}$, and $\rho \mapsto \rho + \langle \gamma, v
  \rangle$ is a bijection of $\Z_q$ (\rocq{com\_dlog},
  \rocq{zk\_of\_bijection}).
  \item Monad laws rearrange the real distribution, which draws all
  randomness at once, into the product form the two steps need
  (\rocq{real\_dist\_proofs}).
\end{enumerate}
The hypothesis on the keys is satisfied by every key in $\langle h_1
\rangle$; the exponents $\gamma$ appear in the proof to exhibit the
bijection and are not used by the simulator.
\end{proof}

\subsection{Cost and assumptions}
Prover: the five commitments and the arguments' announcements, about $(2L
+ 8) n$ exponentiations; verifier about $(2L + 10) n$; the ballot is about
$(2L + 12) n$ field elements plus the $n$ ciphertexts. At $m = 17$: $352$\,KB,
$5.5$\,s to prove and $7.8$\,s to verify, against $6.2$\,MB, $130$\,s and
$68$\,s for the first ballot. Table~\ref{tab:two} compares the two.

\begin{table}[t]
  \centering
  \caption{The two ballots for $m$ candidates, $n = m^2$, $L = \lceil
  \log_2(m+1) \rceil$.}
  \label{tab:two}
  \begin{tabular}{@{}lll@{}}
    \toprule
     & first ballot & batched ballot \\
    \midrule
    ciphertexts & $m + n$ & $n$ \\
    commitments & --- & $5$ (keys of $n$ and $Ln$ generators) \\
    proofs & $m$ membership, $n$ 0/1, $n$ membership & $1$ linear, $3$ zero arguments \\
    size & $O(m^3)$ group elements & $O(n \log m)$ field elements \\
    prover, verifier & $O(m^3)$ exponentiations & $O(n \log m)$ exponentiations \\
    oracle queries & $1$ & $10$ \\
    soundness & exact, $2$ transcripts & statistical, $O(n \log m)$ transcripts \\
    group hypotheses & $g$ of order $q$; char.\ $> 4m + 2$ & char.\ $> 3 \cdot 2^L + 2$ \\
    key hypotheses & none & binding of the commitment keys \\
    zero-knowledge & perfect, no hypothesis & perfect, keys in $\langle h_1 \rangle$ \\
    end to end & \rocq{schulze\_end\_to\_end} & \rocq{schulze\_end\_to\_end2} \\
    \bottomrule
  \end{tabular}
\end{table}

\subsection{A simpler batched ballot}\label{sec:ballot2s}
The batched ballot can be simplified further, and the simplification
makes it smaller and faster rather than slower. It rests on one
observation: once every vector is committed under the same key, most
checks can be done by the verifier through the homomorphism instead of
being proved. This version is mechanised in
\rocq{Frontend/SchulzeSimple.v}: the honest ballot is accepted
(\rocq{ballot3\_valid}) and special soundness holds
(\rocq{ballot3\_sound}), in the same rewinding form as for the batched
ballot, and it is zero-knowledge (\rocq{ballot3\_zk}); the full protocol
is in Section~\ref{sec:ballot2sd}, together with the end-to-end theorem
(\rocq{schulze\_end\_to\_end3}), and Section~\ref{sec:steps3} walks
through it step by step. It is extracted and benchmarked like the
batched ballot, and it was also implemented in the toy group of
Section~\ref{sec:walk} to check the cheating cases. It has four changes.

The first change commits to $s$ under a derived key, so that $C_s$ and
the field row for $\gamma_1$ disappear. From the key $g_1, \ldots, g_n$
anyone can compute a key for positions,
\[
  g'_i = \prod_k g_k^{M_{ki}} = \prod_{j} g_{(j,i)}\, g_{(i,j)}^{-1},
  \qquad i < m,
\]
which is public, like $M$. A commitment to $s$ under $\mathbf{g}'$ is
automatically a commitment to $Ms$ under $\mathbf{g}_1$:
\[
  \com_{\mathbf{g}'}(s; \rho) = h_1^{\rho} \prod_i {g'_i}^{s_i} = h_1^{\rho}
  \prod_k g_k^{(Ms)_k} = \com_{\mathbf{g}_1}(Ms; \rho).
\]
So the voter posts the single commitment $C_\delta = \com_{\mathbf{g}'}(s;
\rho_\delta)$, and $\delta = Ms$ holds by construction. The linear
argument keeps the one row $C_\delta = \com_{\mathbf{g}'}(s;
\rho_\delta)$, so that the extractor still recovers the positions. The
derived key has the known relation $\prod_i g'_i = 1$, because adding a
constant to every position leaves $\delta$ unchanged; this is harmless,
since positions are only read relative to $s_0$ anyway. In the toy group,
$\mathbf{g}' = (32, 8, 9)$ and indeed $32 \cdot 8 \cdot 9 = 1$ modulo
$47$. (That $g'_A = h_1$ there is a coincidence of the tiny group; with
keys hashed into a large group it does not happen.)

The second change commits to the bit vectors $X^{(t)}$ separately, under
the short key $\mathbf{g}_1$, instead of one commitment to $X$ under the
long key $\mathbf{g}_2$. This costs $L - 1$ extra group elements, and the
key $\mathbf{g}_2$ disappears. More importantly, every committed vector
is now under the same key, which the next two changes need.

The third change checks the definition of $v$ through the homomorphism,
so that the field row for $\gamma_2$ disappears. Since $v = -\delta + 2p -
b$ and $v = \sum_t 2^t X^{(t)}$, the verifier multiplies commitments it
already has:
\[
  C_0 = C_\delta \cdot C_b \cdot C_p^{-2} \cdot \prod_t C_{X^{(t)}}^{\,2^t}
      = \com\big(\delta + b - 2p + \textstyle\sum_t 2^t X^{(t)};\; \rho^*\big),
  \qquad \rho^* = \rho_\delta + \rho_b - 2\rho_p + \textstyle\sum_t 2^t \rho_{X^{(t)}} .
\]
For an honest voter the committed vector is exactly zero, so $C_0 =
h_1^{\rho^*}$, and the voter proves knowledge of $\rho^*$ with a Schnorr
row. By binding, this gives the relation entry by entry, exactly, rather
than in a random direction $\gamma_2$. The linear argument no longer
needs $\delta$, $p$, or $X$ in its witness.

The fourth change merges the three zero arguments into one. Every
quadratic check now runs over vectors of length $n$ under one key: the
product $b \circ \delta - p = 0$ (two rows, $b$ against $\delta$ and $p$
against $-\mathbf{1}$), the $L$ bit checks $X^{(t)} \circ (X^{(t)} -
\mathbf{1}) = 0$, and the bit check $b \circ (b - \mathbf{1}) = 0$. With
a new challenge $z$, the voter weights the product relation by $z$, the
$t$-th bit relation by $z^{2+t}$, and the bit relation on $b$ by
$z^{2+L}$, and proves that the weighted sum is zero with a single zero
argument of $L + 3$ rows. The weights are applied to the commitments by
the homomorphism, $C_a \mapsto C_a^{z^i}$. They keep errors in different
relations from cancelling, in the same way that the powers of $y$ keep
coordinates apart: the weighted sum is a polynomial in $z$ whose
coefficients are the individual relations, so if it vanishes for $L + 3$
distinct values of $z$, every relation vanishes. The response of a zero
argument has about $2n$ field elements whatever its number of rows, so
this is the largest saving.

The simplified voter therefore posts the $n$ ciphertexts, the $L + 3$
commitments $C_\delta, C_b, C_p, C_{X^{(0)}}, \ldots, C_{X^{(L-1)}}$, one
linear argument, and one zero argument. The linear argument has the
witness $(s, b, \rho_\delta, \rho_b, R, \rho^*)$ of $m + n + 4$ field
elements and five rows:
\[
  \com_{\mathbf{g}'}(s; \rho_\delta) = C_\delta,\quad
  \com(b; \rho_b) = C_b,\quad
  g^R = C_{\beta,1},\quad
  g^{\langle \beta, b \rangle} h^R = C_{\beta,2},\quad
  h_1^{\rho^*} = C_0 .
\]
The challenges are $\beta$ and $y, z$ from the ciphertexts and
commitments, then $c$ and $x$ from the announcements: five, against ten.

Soundness follows the same case analysis as before. The linear argument
gives $s$, the committed $b$, and the link to the ciphertexts; the zero
argument gives openings of all the commitments satisfying the product
and bit relations; binding then gives $\delta = Ms$ from the derived key
and $v = \sum_t 2^t X^{(t)}$ from $C_0$. Two new lemmas would be needed in
Rocq: the derived-key argument, which uses only the binding of
$\mathbf{g}_1$, and the separation of relations by $z$, a polynomial
argument of the same kind as the existing one for $y$. The existing
$r$-row zero argument is reused as it stands. The soundness error grows by
about $(L + 3)/q$ for $z$, which is negligible at $256$-bit $q$.

In the toy implementation the honest ballot for $A > B > C$ is accepted in
every one of $200$ runs with random challenges, and $C_0 = h_1^{\rho^*}$
holds as claimed. The cheating ballots of Table~\ref{tab:necessary} are
rejected: $b_{AA} = -1$ on a tied cell, a fake $p_{BA} = 0$, and an
encrypted $b_{CA} = 1$ behind an honest commitment are rejected in all
$200$ runs, the first two by the zero argument and the last by the
linear argument. The non-bit $X^{(0)}$ is rejected in $191$ of $200$ runs;
the other $9$ are exactly the runs where the random $y$ is a root of the
cheat's error polynomial, which happens only because the toy field has
$23$ elements.

\begin{table}[ht]
  \centering
  \caption{The batched ballot and its simplification at $m = 17$, $n =
  289$, $L = 5$, at the $2048$-bit parameters with $256$-bit $q$. Sizes
  and times are measured with the extracted programs \rocq{batched.exe}
  and \rocq{simple.exe} in one run on the same machine; the sizes agree
  with the counts of Section~\ref{sec:ballot2sd}, Step~10.}
  \label{tab:simpler}
  \begin{tabular}{@{}lll@{}}
    \toprule
     & batched ballot & simplified \\
    \midrule
    commitment keys & two: $n$ and $Ln$ generators & one: $n$ generators, and $\mathbf{g}'$ \\
    commitments & $5$ & $L + 3 = 8$ \\
    linear witness & $m + (3 + L)n + 6 = 2{,}335$ & $m + n + 4 = 310$ \\
    zero arguments & $3$, responses $(2L + 4)n$ in all & $1$, response about $2n$ \\
    challenges & $10$ & $5$ \\
    proof part of the ballot & $204$\,KB & $36$\,KB \\
    whole ballot & $352$\,KB & $180$\,KB \\
    prove / verify & $4.2$\,s / $5.7$\,s & $0.67$\,s / $1.5$\,s \\
    \bottomrule
  \end{tabular}
\end{table}

The thirteen Gentoo council elections of Section~\ref{sec:problem} run
through this ballot (\rocq{gentoo3.exe}, one process per election, voters
and verifier together) take, per ballot, $0.8$\,s to encrypt and prove at
$m = 10$ and $2.9$\,s at $m = 17$ for the three-pass client with the
voter-bound oracle, and $0.5$\,s to $1.6$\,s to verify; every ballot is
accepted, every replayed ballot rejected, and the certified winners are
the same. The bulletin boards are $8$\,MB ($2023$, $N = 50$) to $43$\,MB
($2013$, $N = 96$), against $175$\,MB to $1.5$\,GB with the first ballot,
whose per-ballot times were $26$ to $132$\,s to prove and $12$ to $57$\,s
to verify. The whole ballot is dominated by its $289$ ciphertexts, about
$148$\,KB, so compressing the proofs further, for instance to the
logarithmic size of Bulletproofs~\cite{bulletproofs}, would add little.
What the simplification gives up is the one-to-one match between the four
facts of the validity statement and the four arguments, which makes the
current version easier to teach. Section~\ref{sec:ballot2sd} gives the
protocol in full.

\subsection{The simpler ballot in full}\label{sec:ballot2sd}
This subsection gives the simpler ballot as a complete protocol, with
every number for $A > B > C$ in the toy group, its soundness and
zero-knowledge theorems, and the size calculation. The definitions, the
honest voter, completeness, soundness, and zero-knowledge are in
\rocq{Frontend/SchulzeSimple.v}; the size count is on paper. The numbers
come from a toy implementation that
follows the definitions of \rocq{zero\_prove} and \rocq{zero\_verify}, and
the keys are those of Section~\ref{sec:walk}: $h_1 = 32$ and $\mathbf{g}_1
= (34, 8, 27, 14, 36, 3, 17, 42, 12)$, with $m = 3$, $n = 9$, $L = 2$.

\subsubsection*{Step 1: the public keys}
Anyone computes the position key $\mathbf{g}'$ from $\mathbf{g}_1$ and the
rule for $M$. Row $(i, j)$ of $M$ has $+1$ in column $j$ and $-1$ in column
$i$, so $g'_i$ collects $g_{(j, i)}$ for every pair that ends at $i$, and
$g_{(i, j)}^{-1}$ for every pair that starts at $i$. The diagonal
generators cancel. For $A, B, C$:
\begin{align*}
  g'_A &= g_{BA}\, g_{CA}\, g_{AB}^{-1}\, g_{AC}^{-1} = 14 \cdot 17 \cdot 8^{-1} \cdot 27^{-1} = 32,\\
  g'_B &= g_{AB}\, g_{CB}\, g_{BA}^{-1}\, g_{BC}^{-1} = 8 \cdot 42 \cdot 14^{-1} \cdot 3^{-1} = 8,\\
  g'_C &= g_{AC}\, g_{BC}\, g_{CA}^{-1}\, g_{CB}^{-1} = 27 \cdot 3 \cdot 17^{-1} \cdot 42^{-1} = 9 .
\end{align*}
No new setup is needed, and no new assumption: binding of $(h_1,
\mathbf{g}_1)$ is the only one the protocol uses.

\subsubsection*{Step 2: what the voter computes and posts}
The voter computes $\delta = Ms$, $b$, $p = b \circ \delta$, $v = -\delta
+ 2p - b$, and the bit vectors $X^{(t)}$, exactly as in the table of
Section~\ref{sec:steps}. The voter posts the $n$ ciphertexts $c_k =
\Enc(b_k; r_k)$ and $L + 3$ commitments:
\[
  C_\delta = \com_{\mathbf{g}'}(s; \rho_\delta), \quad
  C_b = \com(b; \rho_b), \quad
  C_p = \com(p; \rho_p), \quad
  C_{X^{(t)}} = \com(X^{(t)}; \rho_{X^{(t)}}) \ \ (t < L).
\]
Everything but $C_\delta$ is under $\mathbf{g}_1$, and $C_\delta$ is too,
because $\com_{\mathbf{g}'}(s; \rho) = \com_{\mathbf{g}_1}(Ms; \rho)$.

\ex{With $\rho_\delta = 14$, $\rho_b = 3$, $\rho_p = 20$, $\rho_{X^{(0)}}
= 9$, $\rho_{X^{(1)}} = 5$: $C_\delta = 32^{14} \cdot 32^{0} \cdot 8^{1}
\cdot 9^{2} = 27$, which is the $C_\delta = 27$ of
Section~\ref{sec:walk}, computed there from $\delta$ under
$\mathbf{g}_1$. The voter never handles $\delta$ in this computation.
Also $C_b = 4$, $C_p = 17$, $C_{X^{(0)}} = 42$, and $C_{X^{(1)}} = 21$.
The ciphertexts are those of Section~\ref{sec:walk}, with randomness $r =
(3, 6, 11, 2, 13, 17, 20, 8, 10)$.}

\subsubsection*{Step 3: the first challenges and the derived values}
After the ciphertexts and commitments are fixed, the verifier draws $\beta
\in \Z_q^{n}$ and non-zero scalars $y, z$. Both sides then compute, by the
homomorphisms,
\[
  C_\beta = \prod_k c_k^{\beta_k}, \qquad
  C_0 = C_\delta \cdot C_b \cdot C_p^{-2} \cdot \prod_{t < L} C_{X^{(t)}}^{\,2^{t}} .
\]
The voter also knows the openings, $R = \langle \beta, r \rangle$ and
$\rho^* = \rho_\delta + \rho_b - 2\rho_p + \sum_t 2^{t} \rho_{X^{(t)}}$.
Because the committed vector $\delta + b - 2p + \sum_t 2^{t} X^{(t)}$ is
zero for an honest voter, $C_0 = h_1^{\rho^*}$.

\ex{$\beta = (2, 5, 1, 7, 3, 4, 6, 8, 9)$, $y = 4$, $z = 3$. Then $C_\beta
= (32, 7)$ and $R = 5$, as in Section~\ref{sec:walk}. Also $C_0 = 27 \cdot
4 \cdot 17^{-2} \cdot 42 \cdot 21^{2} = 8$, and $\rho^* = 14 + 3 - 40 + 9
+ 10 \equiv 19$, with $h_1^{19} = 32^{19} = 8$.}

\subsubsection*{Step 4: the linear argument}
The witness is $w = (s, b, \rho_\delta, \rho_b, R, \rho^*)$ of length $W'
= m + n + 4$, and the linear map $\varphi' : \Z_q^{W'} \to \G^{5}$ is
\[
  \varphi'(w) = \big(\com_{\mathbf{g}'}(s; \rho_\delta),\;
  \com(b; \rho_b),\; g^{R},\; g^{\langle \beta, b \rangle} h^{R},\;
  h_1^{\rho^*}\big),
\]
with target $T' = (C_\delta, C_b, C_{\beta,1}, C_{\beta,2}, C_0)$. Every
component is linear in $w$, so Section~\ref{sec:linsigma} applies: the voter
draws $a \in \Z_q^{W'}$, announces $A = \varphi'(a)$, receives $c$, and
answers $u = a + c\,w$. The verifier checks $\varphi'(u) = A \cdot
T'^{\,c}$ componentwise. The rows do three jobs. The first shows that
$C_\delta$ is in the image of $M$, the second and the next two tie the
ciphertexts to $C_b$, and the fifth shows that $C_0$ commits to the zero
vector. The field rows of the batched ballot are gone: the two relations
they checked in a random direction now hold entry by entry, by binding.

\ex{The witness, the announcement randomness, and the response at $c =
5$, each split as $(s;\, b;\, \rho_\delta, \rho_b, R, \rho^*)$, are
\begin{align*}
  w &= (0, 1, 2;\ 0, 1, 1, 0, 0, 1, 0, 0, 0;\ 14, 3, 5, 19),\\
  a &= (2, 4, 1;\ 5, 1, 0, 3, 2, 2, 4, 1, 6;\ 7, 8, 3, 10),\\
  u &= (2, 9, 11;\ 5, 6, 5, 3, 2, 7, 4, 1, 6;\ 8, 0, 5, 13),
\end{align*}
and the announcement is $A = \varphi'(a) = (8, 8, 8, 14, 16)$. The verifier's five checks, where
$\langle \beta, u_b \rangle = 2$:
\begin{align*}
  \com_{\mathbf{g}'}((2, 9, 11); 8) &= 7 = 8 \cdot 27^{5}, &
  \com(u_b; 0) &= 14 = 8 \cdot 4^{5}, \\
  g^{5} &= 32 = 8 \cdot 32^{5}, &
  g^{2} h^{5} &= 16 = 14 \cdot 7^{5}, \\
  h_1^{13} &= 3 = 16 \cdot 8^{5}.
\end{align*}}

\subsubsection*{Step 5: one zero argument for all the products}
The quadratic relations are combined by powers of $z$ on the left-hand
vectors. With $r = L + 3$ rows, the statement is
$\sum_{i=1}^{r} a_i \star_y b_i = 0$ for
\begin{center}
\begin{tabular}{@{}llll@{}}
\toprule
row & left $a_i$ & right $b_i$ & commitments, left and right \\
\midrule
$1$ & $z\, b$ & $\delta$ & $C_b^{\,z}$, $C_\delta$ \\
$2$ & $z\, p$ & $-\mathbf{1}$ & $C_p^{\,z}$, $\com(-\mathbf{1}; 0)$ \\
$3 + t$ & $z^{2+t} X^{(t)}$ & $X^{(t)} - \mathbf{1}$ & $C_{X^{(t)}}^{\,z^{2+t}}$, $C_{X^{(t)}} \cdot \com(-\mathbf{1}; 0)$ \\
$3 + L$ & $z^{2+L} b$ & $b - \mathbf{1}$ & $C_b^{\,z^{2+L}}$, $C_b \cdot \com(-\mathbf{1}; 0)$ \\
\bottomrule
\end{tabular}
\end{center}
Expanding the sum as a polynomial in $z$ shows what it says:
\[
  z\,\big(b \star_y \delta - p \star_y \mathbf{1}\big)
  + \sum_{t < L} z^{2 + t}\, X^{(t)} \star_y (X^{(t)} - \mathbf{1})
  + z^{2 + L}\, b \star_y (b - \mathbf{1}) = 0 .
\]
The verifier computes all the commitments in the table from the posted
ones, so no new commitment is posted. Rows $1$ and $2$ need the same
weight, because only their sum vanishes; every other row vanishes by
itself.

The argument is the zero argument of Section~\ref{sec:zero} with $r$ rows,
in the form the development uses. The voter draws blinding vectors $a_0,
b_{\mathrm{m}} \in \Z_q^{n}$ with randomness $r_0, s_{\mathrm{m}}$, and
defines
\[
  A(x) = a_0 + \sum_{i=1}^{r} x^{i} a_i, \qquad
  B(x) = b_{\mathrm{m}} + \sum_{i=1}^{r} x^{r + 1 - i} b_i .
\]
In the product $A(x) \star_y B(x) = \sum_{k=0}^{2r} d_k x^{k}$, the pair
$a_i, b_i$ always lands at $x^{r+1}$, so the coefficient $d_{r+1}$ is the
claimed sum, which is zero. The voter announces $\com(a_0; r_0)$,
$\com(b_{\mathrm{m}}; s_{\mathrm{m}})$, and $D_k = \com(d_k; t_k)$ for $k
\le 2r$, over the first generator of the key, with $t_{r+1} = 0$, so that
$D_{r+1} = 1$ is visibly a commitment to zero. After the challenge $x$ the
voter answers with $A(x)$, $B(x)$, and the randomness $r_x$, $s_x$, $t_x$
obtained by evaluating the same polynomials on the randomness. The
verifier checks $D_{r+1} = 1$ and
\begin{gather*}
  \com(a_0; r_0) \prod_{i} (C_{a_i})^{x^{i}} = \com(A(x); r_x), \qquad
  \com(b_{\mathrm{m}}; s_{\mathrm{m}}) \prod_{i} (C_{b_i})^{x^{r+1-i}} = \com(B(x); s_x),\\
  \prod_{k} D_k^{x^{k}} = \com(A(x) \star_y B(x); t_x).
\end{gather*}

\ex{$y = 4$, $z = 3$, $r = 5$, $x = 6$. The weights are $z = 3$, $z^2 = 9$,
$z^3 = 4$, $z^4 = 12$, so the rows are
\begin{align*}
  a_1 &= (0,3,3,0,0,3,0,0,0), & b_1 &= (0,1,2,22,0,1,21,22,0),\\
  a_2 &= (0,3,6,0,0,3,0,0,0), & b_2 &= (22,22,22,22,22,22,22,22,22),\\
  a_3 &= (0,0,9,9,0,0,0,9,0), & b_3 &= (22,22,0,0,22,22,22,0,22),\\
  a_4 &= (0,0,0,0,0,0,4,0,0), & b_4 &= (22,22,22,22,22,22,0,22,22),\\
  a_5 &= (0,12,12,0,0,12,0,0,0), & b_5 &= (22,0,0,22,22,0,22,22,22).
\end{align*}
The five values $a_i \star_4 b_i$ are $1, 22, 0, 0, 0$, which sum to $0$
modulo $23$. Rows $1$ and $2$ cancel each other, and the three bit rows
vanish alone. With $\com(-\mathbf{1}; 0) = 42$, the statement's commitments
are
\[
  (C_{a_i}) = (17, 25, 7, 42, 2), \qquad (C_{b_i}) = (27, 42, 25, 36, 27).
\]
The voter announces, with $a_0 = (5,2,7,1,9,4,3,8,6)$, $r_0 = 11$,
$b_{\mathrm{m}} = (3,6,1,8,2,5,7,4,9)$, $s_{\mathrm{m}} = 13$, and $t =
(17,19,4,8,2,10,0,1,12,21,15)$:
\begin{align*}
  &\com(a_0; 11) = 37, \qquad \com(b_{\mathrm{m}}; 13) = 36,\\
  &d = (22,21,2,2,4,20,\mathbf{0},19,14,10,4),\\
  &D = (42,7,27,21,9,17,\mathbf{1},1,6,28,27).
\end{align*}
(The second $1$ in $D$ is a coincidence of the tiny group, where $g_1$ and
$h_1$ have small relations.) The response is
\begin{align*}
  A(6) &= (5,14,1,13,9,16,12,20,6), & r_x &= 12,\\
  B(6) &= (13,1,7,2,12,0,3,21,19), & s_x &= 21, \qquad t_x = 9.
\end{align*}
The verifier checks $D_6 = 1$, then that both sides of the first equation
are $4$, both sides of the second are $21$, and both sides of the third are
$12$, with $A(6) \star_4 B(6) = 13$ and $\com(13; 9) = 12$.}

\subsubsection*{Step 6: the verifier and the Fiat--Shamir order}
The verifier checks that there are $n$ ciphertexts, computes $C_\beta$,
$C_0$, and the statement of the zero argument, checks the five equations
of the linear argument, and checks the zero argument. The five challenges
are hashed in this order:
\begin{itemize}[nosep]
  \item $\beta$ (one query per coordinate), $y$, and $z$, each with its own
  tag, from the ciphertexts and the $L + 3$ commitments;
  \item $c$ from those and the linear announcement $A$;
  \item $x$ from those and the announcement of the zero argument.
\end{itemize}
The honest voter computes the ballot in three passes, as before. The first
fixes the ciphertexts and commitments, hence $\beta$, $y$, $z$. The second
computes $A$, which depends on $\beta$, and the zero argument's
announcement, whose $d$ depends on $y$ and $z$, hence $c$ and $x$. The
third computes the responses.

\subsubsection*{Step 7: soundness}
The statement below is the mechanised theorem \rocq{ballot3\_sound}, and
the proof follows the mechanised one. It has the same shape as
Theorem~\ref{thm:b2sound}, and it assumes the same bound on the
characteristic of the field, above $3 \cdot 2^{L} + 2$.

\begin{proposition}[\rocq{ballot3\_sound}]\label{prop:v3sound}
Suppose a simpler ballot is accepted by the linear argument at a base
$\beta$ with two challenges $c \ne c'$ on one announcement, and for every
coordinate $k$ at a vector differing from $\beta$ only in $\beta_k$, again
with two challenges; and by the zero argument at $n$ distinct non-zero $y$
and, for each, $L + 2$ distinct non-zero $z$, with $2r + 1 = 2L + 7$ distinct
challenges $x$ on one announcement for each pair $(y, z)$. Then either the
binding of $(h_1, \mathbf{g}_1)$ is broken, or there is $\pos \in \{0,
\ldots, 2^{L+1}\}^m$ such that every $c_{ij}$ decrypts to
$g^{[\pos_i < \pos_j]}$.
\end{proposition}
\begin{proof}
\emph{The linear witness.} The two challenges at the base give, by the
lemma of Section~\ref{sec:linsigma}, values $s$, $b$, $\rho_\delta$,
$\rho_b$, $R$, $\rho^*$ with $C_\delta = \com_{\mathbf{g}'}(s; \rho_\delta)
= \com(Ms; \rho_\delta)$, $C_b = \com(b; \rho_b)$, $C_\beta = \Enc(\langle
\beta, b \rangle; R)$, and $C_0 = h_1^{\rho^*}$.

\emph{The link.} Each $\beta$ differing in coordinate $k$ gives a witness
opening $C_b$; it opens to the same $b$ unless binding is broken. The
division-and-root argument (\rocq{beta\_bits}) then gives $\Dec(c_k) =
g^{b_k}$ for every $k$.

\emph{The zero argument.} At each $(y, z)$, the $2r + 1$ transcripts
extract (\rocq{zero\_soundness}) openings of all the commitments in the
table of Step 5 with $\sum_i a_i \star_y b_i = 0$, or a binding break.
Dividing by $z^{i}$, which is non-zero, turns an opening of $C^{z^{i}}$
into an opening of $C$. By binding, all these openings agree: there are
fixed $\delta'$, $b$, $p$, and $X^{(t)}$ opening $C_\delta$, $C_b$, $C_p$,
and $C_{X^{(t)}}$ at every $(y, z)$. By binding again, the right-hand
openings are $-\mathbf{1}$, $X^{(t)} - \mathbf{1}$, and $b - \mathbf{1}$,
since multiplying by $\com(-\mathbf{1}; 0)^{-1}$ maps them to openings of
$C_{X^{(t)}}$ and $C_b$.

\emph{Separating $z$, then $y$.} For fixed $y$, divide the sum in Step 5
by $z$. This leaves a polynomial in $z$ of degree $L + 1$ with $L + 2$
distinct roots, so every coefficient vanishes: $b \star_y \delta' - p
\star_y \mathbf{1} = 0$, $X^{(t)} \star_y (X^{(t)} - \mathbf{1}) = 0$, and
$b \star_y (b - \mathbf{1}) = 0$. Each of these holds at $n$ distinct
non-zero $y$, so the argument of Section~\ref{sec:zero} gives $p = b \circ
\delta'$, and every entry of $X^{(t)}$ and of $b$ is a bit.

\emph{The derived key.} $C_\delta$ has the openings $(\delta',
\rho_{\delta'})$ and $(Ms, \rho_\delta)$ under $\mathbf{g}_1$. By binding,
$\delta' = Ms$.

\emph{The expansion of $v$.} By the homomorphism, $C_0$ is a commitment to
$\delta' + b - 2p + \sum_t 2^{t} X^{(t)}$. It is also $h_1^{\rho^*} =
\com(\mathbf{0}; \rho^*)$. By binding, $\sum_t 2^{t} X^{(t)} = -\delta' + 2p
- b$.

\emph{Decoding.} These are the relations reached at Step 6 of the proof of
Theorem~\ref{thm:b2sound}, and the same decoding gives $\pos$
(\rocq{decode\_positions}). With the link, $c_{ij}$ decrypts to
$g^{[\pos_i < \pos_j]}$.
\end{proof}

The number $L + 2$ is one less than the $L + 3$ rows because the polynomial
in $z$ has no constant term. Counting the bad events as usual, a cheating
voter's probability of success is at most about
\[
  \frac{1}{q} \ (\text{for } \beta) + \frac{n}{q} \ (\text{for } y)
  + \frac{L + 1}{q} \ (\text{for } z) + \frac{1}{q} \ (\text{for } c) +
  \frac{2L + 6}{q} \ (\text{for } x),
\]
which is negligible at $256$-bit $q$. This count is an estimate, not a
mechanised bound; the development's theorems, here as elsewhere, stop at
the rewinding form above.

\subsubsection*{Step 8: zero-knowledge}
This is the mechanised theorem \rocq{ballot3\_zk}, in the same
permutation form as \rocq{ballot2\_zk}. The simulator follows that proof, with commitment keys
of the form $g_{1,j} = h_1^{\gamma_j}$. It draws $L + 3$ uniform exponents
and sets the commitments to the corresponding powers of $h_1$. It computes
$C_\beta$, $C_0$, and the statement of the zero argument from them, and it
runs the simulators of the linear argument and of the zero argument.
Commitments with fresh randomness are uniform in $\G$ whatever they commit
to, so these are distributed as in the real ballot. The equation $C_0 =
h_1^{\rho^*}$ does not distinguish the two, because $h_1$ generates $\G$,
so every group element has this form. The derived key causes no difficulty:
the relation $\prod_i g'_i = 1$ only means that $C_\delta$ has many openings
under $\mathbf{g}'$, and the factor $h_1^{\rho_\delta}$ still makes
$C_\delta$ uniform.
The ciphertexts, as in both other ballots, are covered by IND-CPA.

\subsubsection*{Step 9: testing on cheating ballots}
Running this protocol on random challenges in the toy group, with the keys
above, gives the following. The cheater here is the honest algorithm run
on the wrong vectors, with $D_{r+1}$ set to $1$ regardless; a cleverer
cheater can also try to guess $x$, which is the last term of the count
above.
\begin{center}
\begin{tabularx}{\textwidth}{@{}>{\raggedright\arraybackslash}Xlr@{}}
\toprule
ballot & runs and acceptances & where it fails \\
\midrule
honest & $500$, all accepted & --- \\
$b_{AA} = -1$, with $X^{(0)}_{AA} = 1$ & $2000$, none accepted & zero argument \\
$b_{BA} = 1$, fake $p_{BA} = 0$, $v_{BA} = 0$ & $2000$, none accepted & zero argument \\
$X^{(0)}_{CA} = 2$, $X^{(1)}_{CA} = 0$ & $2000$, none accepted & zero argument \\
$\delta_{BA} = 1$ not from $s$, $b_{BA} = p_{BA} = 1$ & $4000$, $153$ accepted & linear, row $C_\delta$ \\
$X^{(0)}_{AA} = 1$, which breaks $C_0$ & $4000$, $147$ accepted & linear, row $C_0$ \\
$b_{CA} = 1$ encrypted, honest $C_b$ & $4000$, $322$ accepted & linear, link rows \\
\bottomrule
\end{tabularx}
\end{center}
The three zero-argument cheats each leave one wrong entry, so their error
is a single term $e\, y^{j} z^{i}$, which has no non-zero root in $y$ or
$z$, and this cheater always fails. The linear cheats are accepted exactly
when $c = 0$, and in the last row also when $\beta_{CA} = 0$: these are the
$1/23$ events of the tiny field, and the script confirmed in all $12{,}000$
runs that acceptance happens exactly in these cases.

\subsubsection*{Step 10: the size, term by term}
Let $\lvert \G \rvert$ and $\lvert \Z_q \rvert$ be the sizes of a group and
a field element, $256$ and $32$ bytes at the parameters of
Table~\ref{tab:simpler}. The whole ballot has
\begin{center}
\begin{tabular}{@{}lll@{}}
\toprule
part & group elements & field elements \\
\midrule
ciphertexts & $2n$ & \\
commitments & $L + 3$ & \\
linear argument & $5$ & $m + n + 4$ \\
zero argument & $2r + 3 = 2L + 9$ & $2n + 3$ \\
\bottomrule
\end{tabular}
\end{center}
The element $D_{r+1} = 1$ is counted in the zero argument but could be
omitted. At $m = 17$, $n = 289$, $L = 5$, the proof part is $3L + 17 = 32$
group and $m + 3n + 7 = 891$ field elements, that is, $8.2 + 28.5 \approx
37$\,KB. The ciphertexts add $578 \cdot 256$ bytes, about $148$\,KB, giving
about $185$\,KB. The extracted ballot (\rocq{simple.exe}, with the keys
hashed to the group as for the batched ballot) has exactly $610$ group and
$891$ field elements, $180$\,KB, and takes $0.67$\,s to prove and
$1.5$\,s to verify at $m = 17$; the batched ballot, run in the same
session on the same machine, takes $4.2$\,s and $5.7$\,s. The small gap
between $185$ and $180$\,KB is the rounding of the per-element sizes.

\subsubsection*{Step 11: what the mechanisation reused}
The mechanisation, \rocq{Frontend/SchulzeSimple.v}, is about $1{,}300$
lines with no axioms. It reuses the zero argument
(\rocq{zero\_soundness}), the linear argument (\rocq{lin\_soundness}),
the link (\rocq{beta\_bits}), and the decoding
(\rocq{decode\_positions}) as they stand. Three things are new: the
definition of the ballot with its $\varphi'$ (\rocq{phi3}) and the rows
of the zero argument as functions of the row index (\rocq{rowA},
\rocq{rowB}); the derived-key lemma $\com_{\mathbf{g}'}(s; \rho) =
\com_{\mathbf{g}_1}(Ms; \rho)$ (\rocq{com\_gd}), which is a reordering of
a double product (\rocq{com\_lincomb}); and the separation of the rows by
$z$ (\rocq{weighted\_sum\_eval}, \rocq{coeffs\_vanish}), the same
Vandermonde argument as for $y$ (\rocq{poly\_vanish\_mul}). The openings
of $C_p$ and the $C_{X^{(t)}}$ are taken from one accepted zero argument
(\rocq{zero\_open\_ref}) and shown by binding to be the ones every other
accepted $(y, z)$ opens (\rocq{zero\_open\_at}). The zero-knowledge proof
(\rocq{ballot3\_zk}) is the batched one with $L + 3$ commitment
randomnesses in the bijection and one zero argument instead of three;
$C_\delta$ is rewritten as a commitment under $\mathbf{g}_1$ by
\rocq{com\_gd}, so the derived key needs no discrete logarithms of its
own.

\section{A ballot from start to finish: $A > B > C$}\label{sec:walk}
This section takes one ballot, $A > B > C$, and writes down every number
of both protocols in the toy group ($p = 47$, $q = 23$, $g = 2$, $x = 5$,
$h = 32$). All randomness is chosen by hand so that the reader can redo
each equation; a real client draws it uniformly. The ballot is then
aggregated with ballot 1 of the running example and the two-ballot
election is certified.

\subsection{Positions and bits}
$A > B > C$ is the position vector $s = (s_A, s_B, s_C) = (0, 1, 2)$. The
six off-diagonal cells, with $\delta_{ij} = s_j - s_i$ and $b_{ij} =
[\delta_{ij} \ge 1]$:
\begin{center}
\begin{tabular}{@{}lrr@{\qquad}lrr@{}}
\toprule
cell & $\delta$ & $b$ & cell & $\delta$ & $b$ \\
\midrule
$(A,B)$ & $1$ & $1$ & $(B,A)$ & $-1$ & $0$ \\
$(A,C)$ & $2$ & $1$ & $(C,A)$ & $-2$ & $0$ \\
$(B,C)$ & $1$ & $1$ & $(C,B)$ & $-1$ & $0$ \\
\bottomrule
\end{tabular}
\end{center}
The ballot contributes one to $D_{AB}$, $D_{AC}$ and $D_{BC}$ and nothing
else, as a ranking $A > B > C$ should.

\subsection{The first ballot}

\paragraph{Ciphertexts} The voter encrypts the positions with randomness
$r_A = 4$, $r_B = 9$, $r_C = 15$:
\[
  \Enc(0; 4) = (16, 6),\qquad \Enc(1; 9) = (42, 1),\qquad \Enc(2; 15) = (9, 21),
\]
and the bits with randomness $6, 11, 2, 17, 20, 8$ in the order of
Table~\ref{tab:walk1}. The verifier derives $\Enc(\delta_{ij}) =
\Enc(s_j) / \Enc(s_i)$ and forms the shifted ciphertext $\Enc(b_{ij})^7
\cdot \Enc(\delta_{ij})$, whose plaintext is $t = 7 b_{ij} + \delta_{ij}$.
For instance, for $(C, A)$: $\Enc(\delta_{CA}) = (16, 6) / (9, 21) = (16
\cdot 9^{-1}, 6 \cdot 21^{-1}) = (7, 7)$, which is $\Enc(-2; 4 - 15) =
\Enc(-2; 12)$; and $(6, 21)^7 \cdot (7, 7) = (28, 24) = \Enc(-2; 7 \cdot
20 + 12) = \Enc(-2; 14)$.

\begin{table}[t]
\centering
\caption{The first ballot for $A > B > C$: the six pair cells. The
randomness of the shifted ciphertext is $7 r_b + r_j - r_i$.}
\label{tab:walk1}
\begin{tabular}{@{}lrrcccr@{}}
\toprule
cell & $\delta$ & $b$ & $\Enc(b; r_b)$ & $\Enc(\delta)$ derived & shifted & $t$ \\
\midrule
$(A,B)$ & $1$ & $1$ & $(17,21) = \Enc(1;6)$ & $(32,8) = \Enc(1;5)$ & $(2,14)$ & $8$ \\
$(A,C)$ & $2$ & $1$ & $(27,37) = \Enc(1;11)$ & $(27,27) = \Enc(2;11)$ & $(3,7)$ & $9$ \\
$(B,A)$ & $-1$ & $0$ & $(4,37) = \Enc(0;2)$ & $(25,6) = \Enc(-1;18)$ & $(42,12)$ & $-1$ \\
$(B,C)$ & $1$ & $1$ & $(36,36) = \Enc(1;17)$ & $(17,21) = \Enc(1;6)$ & $(37,7)$ & $8$ \\
$(C,A)$ & $-2$ & $0$ & $(6,21) = \Enc(0;20)$ & $(7,7) = \Enc(-2;12)$ & $(28,24)$ & $-2$ \\
$(C,B)$ & $-1$ & $0$ & $(21,36) = \Enc(0;8)$ & $(36,9) = \Enc(-1;17)$ & $(16,3)$ & $-1$ \\
\bottomrule
\end{tabular}
\end{table}

Every $t$ is in $\{-3, \ldots, 0\} \cup \{8, 9, 10\}$, as it must be for
an honest ballot. All proofs below use the single challenge $c = 17$; in
the deployed ballot $c$ is the hash of the credential, all ciphertexts and
all announcements.

\paragraph{A position proof in full} Candidate $A$: $C = \Enc(0; 4) =
(16, 6)$, messages $\{0, 1, 2, 3\}$, true branch $k = 0$. The voter
simulates branches $1, 2, 3$ by choosing $(c_k, z_k) = (3, 7), (8, 19),
(12, 2)$ and computing $A_k = (g^{z_k} C_1^{-c_k},\; h^{z_k} (C_2 /
g^{k})^{-c_k})$:
\[
  A_1 = (25, 2),\quad A_2 = (37, 6),\quad A_3 = (1, 14);
\]
and computes the true branch honestly with $a_0 = 13$: $A_0 = (g^{13},
h^{13}) = (14, 3)$. After the challenge $c = 17$: $c_0 = 17 - (3 + 8 +
12) = -6 \equiv 17$, $z_0 = 13 + 17 \cdot 4 = 81 \equiv 12$. The verifier
checks $c_0 + c_1 + c_2 + c_3 \equiv 17$ and, for every branch, $g^{z_k} =
A_{k,1} C_1^{c_k}$ and $h^{z_k} = A_{k,2} (C_2 / g^k)^{c_k}$; for the true
branch: $2^{12} = 7 = 14 \cdot 16^{17}$ and $32^{12} = 28 = 3 \cdot
6^{17}$ (all modulo $47$). The simulated branches hold by
construction. Nothing in the transcript says which branch was honest.

\paragraph{A pair cell in full} Cell $(C, A)$: $\Enc(b_{CA}) = (6, 21)$
with $b_{CA} = 0$, randomness $20$.
\emph{0/1 proof}: branch $1$ simulated with $(c_1, z_1) = (5, 9)$, $A_1 =
(2, 37)$; branch $0$ honest with $a_0 = 21$, $A_0 = (12, 14)$; $c_0 = 17 -
5 = 12$, $z_0 = 21 + 12 \cdot 20 = 261 \equiv 8$. \emph{Membership proof}
on the shifted ciphertext $(28, 24) = \Enc(-2; 14)$, messages $\{-3, -2,
-1, 0, 8, 9, 10\}$, true branch $-2$: simulated $(c_k, z_k)$ for $k = -3,
-1, 0, 8, 9, 10$ are $(2,5), (6,8), (1,14), (4,3), (9,20), (7,11)$ with
announcements $(1,12), (18,36), (1,4), (18,32), (42,17), (32,36)$; the
honest branch uses $a = 3$, $A_{-2} = (8, 9)$, $c_{-2} = 17 - (2 + 6 + 1 +
4 + 9 + 7) = -12 \equiv 11$, $z_{-2} = 3 + 11 \cdot 14 = 157 \equiv 19$.
Check of the honest branch: $g^{19} = 3 = 8 \cdot 28^{11}$ and $h^{19}
= 8 = 9 \cdot (24 \cdot g^{2})^{11}$ (the message is $-2$, so $C_2 /
g^{-2} = C_2 \cdot g^{2}$). The cell costs $2 + 7 = 9$ Chaum--Pedersen transcripts.

\paragraph{What a cheater runs into} Suppose the voter wants $b_{CA} = 1$
instead. The shifted plaintext would be $t = 7 + (-2) = 5$, which is in no
branch's message, so no branch can be run honestly. The voter can
simulate all seven branches, but their challenges $c_k$ are then fixed
before $c = H(\ldots)$ is known, and the verifier's check $\sum_k c_k = c$
holds with probability $1/23$ (in a real group, $1/q \approx 2^{-256}$).
This is the special-soundness argument in one sentence: two different
$c$'s with the same announcements force a difference in some branch's
challenge, and that branch reveals the randomness of the ciphertext for
its message.

Concretely, take the cheating row $(B, A)$ of Table~\ref{tab:ballot1}:
the voter encrypts $b_{BA} = 1$ with randomness $2$, $\Enc(1; 2) = (4,
27)$, while the derived $\Enc(\delta_{BA}) = \Enc(-1; 18)$ is what it is.
The shifted ciphertext is $\Enc(1; 2)^7 \cdot \Enc(-1; 18) = \Enc(6; 9) =
(42, 32)$. The voter knows that $6$ is in no branch, but tries the branch
for message $8$ honestly anyway (the branch an honest $b = 1$ cell would
use): announcement $(g^3, h^3) = (8, 9)$, challenge $c = 17$, response $z
= 3 + 17 \cdot 9 \equiv 18$. The first equation holds, $g^{18} = 25 = 8
\cdot 42^{17}$, because the voter does know the randomness $9$. The
second fails: $h^{18} = 12$ but $9 \cdot (32 / g^{8})^{17} = 9 \cdot
6^{17} = 37$, because $C_2 / g^8 = g^{6 - 8} h^{9}$ is not a pure power
of $h$; the equation would need $g^{-2c} = 1$, i.e.\ $c = 0$. Any other
choice fares the same: a branch run honestly fails its second equation,
and a fully simulated set of branches fails the challenge-sum check.

\subsection{The batched ballot}

\paragraph{The vectors} With $n = 9$ cells in the order $AA, AB, AC, BA,
BB, BC, CA, CB, CC$ and $L = 2$:
\begin{align*}
  s &= (0, 1, 2, 0, 0, 0, 0, 0, 0) \quad\text{(padded to length $9$)}, \\
  \delta &= (0, 1, 2,\; -1, 0, 1,\; -2, -1, 0), \\
  b &= (0, 1, 1,\; 0, 0, 1,\; 0, 0, 0), \\
  p = b \circ \delta &= (0, 1, 2,\; 0, 0, 1,\; 0, 0, 0), \\
  v = -\delta + 2p - b &= (0, 0, 1,\; 1, 0, 0,\; 2, 1, 0), \\
  X^{(0)} &= (0, 0, 1,\; 1, 0, 0,\; 0, 1, 0), \qquad X^{(1)} = (0, 0, 0,\; 0, 0, 0,\; 1, 0, 0).
\end{align*}
Check cell $CA$: $v = 2 + 0 - 0 = 2 = 0 + 2 \cdot 1$. Check cell $AC$: $v
= -2 + 4 - 1 = 1$. Every $v$ is in $\{0, 1, 2, 3\}$ and every $X$ entry
is a bit.

\paragraph{Ciphertexts and commitments} The nine bits are encrypted with
randomness $r = (3, 6, 11, 2, 13, 17, 20, 8, 10)$:
\[
  c = \big((8,9), (17,21), (27,37), (4,37), (14,3), (36,36), (6,21), (21,36), (37,16)\big).
\]
The commitment keys are $h_1 = 32$, $\mathbf{g}_1 = (34, 8, 27, 14, 36,
3, 17, 42, 12)$ and $\mathbf{g}_2 = \mathbf{g}_1 \,\|\, (16, 21, 37, 7,
28, 9, 18, 25, 6)$ (all in the subgroup; in a deployment they are hashed
to the group). With randomness $\rho = (\rho_s, \rho_\delta, \rho_b,
\rho_p, \rho_X) = (6, 14, 3, 20, 9)$:
\[
  C_s = 42,\quad C_\delta = 27,\quad C_b = 4,\quad C_p = 17,\quad C_X = 4 .
\]
(That $C_b = C_X$ is a coincidence of the small group; they are
commitments over different keys.) For instance $C_s = 32^6 \cdot 8^{1}
\cdot 27^{2} = 42$.

\paragraph{The random linear combination} The verifier's coefficients,
drawn after the ballot is posted, are $\beta = (2, 5, 1, 7, 3, 4, 6, 8,
9)$. Then $\langle \beta, b \rangle = 5 + 1 + 4 = 10$ and $\langle \beta,
r \rangle = 5$ (mod $23$), and indeed
\[
  C_\beta = \prod_k c_k^{\beta_k} = (32, 7) = \Enc(10; 5).
\]
The voter proves that $C_\beta$ encrypts the inner product of $\beta$ with
the \emph{committed} $b$, with randomness $R = 5$. Had the voter committed
to $b' = b$ with $b'_{CA} = 1$, the inner product would be $10 + 6 = 16
\ne 10$, and the row $g^{\langle \beta, b' \rangle} h^{R}$ could not match
$C_{\beta, 2}$ for any $R$ consistent with $C_{\beta,1} = g^R$.

\paragraph{The field rows} With $\gamma_1 = (1, \ldots, 9)$ and $\gamma_2
= (9, \ldots, 1)$: $\langle \gamma_1, \delta - Ms \rangle = 0$ and
$\langle \gamma_2, X^{(0)} + 2 X^{(1)} + \delta - 2p + b \rangle = 0$, as
the reader can check entry by entry (each entry of the two vectors is
already $0$). If the voter committed to $\delta'$ with $\delta'_{CA} = 1$
instead of $-2$, the first row would be $\langle \gamma_1, \delta' - Ms
\rangle = 7 \cdot 3 = 21 \ne 0$.

\paragraph{The linear argument, three rows of it} The witness is $w = (s,
\delta, b, p, X, \rho_s, \rho_\delta, \rho_b, \rho_p, \rho_X, R)$ of
length $W = 3 + 27 + 18 + 6 = 54$ (the development pads $s$ to $m$
entries, not $n$; the padding here is only for display). The prover draws
$a \in \Z_{23}^{54}$, announces $A = \varphi(a)$, receives $c = 5$, and
answers $z = a + 5 w$. Three rows, with the relevant slices of $a$
chosen as $a_s = (2, 4, 1, 0, \ldots)$, $a_{\rho_s} = 7$, $a_R = 3$,
$a_\delta = (5, 1, 0, 3, 2, 2, 4, 1, 6)$:
\begin{itemize}[nosep]
  \item Row $C_s$: $A_1 = \com(a_s; a_{\rho_s}) = 8$. Response $z_s = a_s +
  5 s = (2, 9, 11, 0, \ldots)$, $z_{\rho_s} = 7 + 5 \cdot 6 = 37 \equiv
  14$. Check: $\com(z_s; z_{\rho_s}) = 4$ and $A_1 \cdot C_s^{5} = 8 \cdot
  42^5 = 4$.
  \item Row $g^R$: $A_6 = g^{3} = 8$, $z_R = 3 + 5 \cdot 5 = 28 \equiv 5$.
  Check: $g^{5} = 32 = A_6 \cdot C_{\beta,1}^{5} = 8 \cdot 32^5$.
  \item Field row $\gamma_1$: $A_8 = \langle \gamma_1, a_\delta - M a_s
  \rangle = 10$. Check: $\langle \gamma_1, z_\delta - M z_s \rangle = 10 =
  A_8 + 5 \cdot 0$.
\end{itemize}
The other six rows are checked the same way; the whole argument is one
announcement of $7$ group and $2$ field elements and one response of
$54$ field elements.

\paragraph{The zero arguments} At $y_1 = 4$ the product relation is
$\sum_{j} 4^{j+1} (b_j \delta_j - p_j) = 0$ because every term is $0$;
likewise the bit relations $\sum_j 4^{j+1} X_j (X_j - 1) = 0$ and $\sum_j
4^{j+1} b_j (b_j - 1) = 0$. The arguments themselves are the library's
Bayer--Groth zero arguments on the commitments $(C_b, C_p)$ against
$(C_\delta, \com(-\mathbf{1}; 0))$, $C_X$ against $C_X \cdot
\com(-\mathbf{1}; 0)$, and $C_b$ against $C_b \cdot \com(-\mathbf{1};
0)$. The bit argument on $b$, at $y_3 = 4$ and $x = 6$, runs as follows
(every number checked by script; the weights $4^j$ for $j = 1, \ldots, 9$
are $4, 16, 18, 3, 12, 2, 8, 9, 13$):
\begin{align*}
  &\text{claim:} && b \star_4 (b - \mathbf{1}) = 0, \quad C_b = 4, \quad
     C_{b - \mathbf{1}} = C_b \cdot \com(-\mathbf{1}; 0) = 27;\\
  &\text{announce:} && a_0 = (5,2,7,1,9,4,3,8,6),\ \com(a_0; 11) = 37;\\
  & && c_0 = (3,6,1,8,2,5,7,4,9),\ \com(c_0; 13) = 36;\\
  & && d_0 = 22,\ D_0 = \com(22; 17) = 42; \quad d_1 = 3,\ D_1 = \com(3; 19) = 2;\\
  &\text{respond:} && A(6) = a_0 + 6b = (5,8,13,1,9,10,3,8,6), \text{ randomness } 11 + 6 \cdot 3 = 6;\\
  & && C(6) = 6(b - \mathbf{1}) + c_0 = (20,6,1,2,19,5,1,21,3), \text{ randomness } 6 \cdot 3 + 13 = 8;\\
  & && t = 17 + 6 \cdot 19 = 16;\\
  &\text{verify:} && \com(A(6); 6) = 24 = 37 \cdot 4^6, \qquad
     \com(C(6); 8) = 28 = 27^6 \cdot 36,\\
  & && A(6) \star_4 C(6) = 17 = 22 + 6 \cdot 3, \qquad \com(17; 16) = 9 = 42 \cdot 2^6 .
\end{align*}
All arithmetic is modulo $23$ for vectors and scalars and modulo $47$ for
group elements. Had the voter set $b_{CA} = 2$, the claim would be $b
\star_4 (b - \mathbf{1}) = 16$: a non-zero coefficient of $x^2$ with no
commitment to absorb it, so the last check fails for all but at most two
values of $x$. The product argument has two rows and runs the same way
with polynomials of degree $2$ in $x$. What matters for soundness is the
polynomial view in $y$: for the cheating
vector $b'$ above, $b'_{CA} \delta_{CA} - p'_{CA}$ is $-2 - p'_{CA}$, so
either $p'_{CA} = -2$ (and then $v'_{CA} = 2 - 4 - 1 = -3 \equiv 20$, not
a sum of two bits, caught by the bit argument on $X$) or the product
relation has a non-zero coefficient at $y^{7}$ and holds for at most $9$
of the $22$ non-zero values of $y$.

\subsection{Aggregation with ballot 1}
Ballot 1 of the running example ($A > B = C$, $s = (0, 1, 1)$) has bits
$(0, 1, 1, 0, 0, 0, 0, 0, 0)$; encrypt them with randomness $1, \ldots,
9$. Multiplying cell by cell with the ciphertexts of $A > B > C$ above
gives the aggregate
\[
  \big((16,6), (21,3), (28,8), (17,34), (25,12), (1,2), (16,6), (18,27), (3,8)\big),
\]
which the trustees decrypt to $(1, 4, 4, 1, 1, 2, 1, 1, 1) = (g^0, g^2,
g^2, g^0, g^0, g^1, g^0, g^0, g^0)$: the margins of the two-ballot
election are $D_{AB} = 2$, $D_{AC} = 2$, $D_{BC} = 1$ and all others
$0$. Nobody needed a key to compute the aggregate; the key was used once,
by the trustees, in private.

\subsection{The certificate for the two-ballot election}
$N = 2$, $K = 2N + 1 = 5$. Defeats: $d_{AB} = 2$ ($2 > 0$), $d_{AC} = 2$,
$d_{BC} = 1$, all others $0$. Strongest paths: $p_{AB} = 2$, $p_{AC} = 2$,
$p_{BC} = 1$, $p_{BA} = p_{CA} = p_{CB} = 0$. Winner $A$: $p_{AB} = 2 \ge
p_{BA} = 0$ and $p_{AC} = 2 \ge p_{CA} = 0$. The trustees publish:
\begin{itemize}[nosep]
  \item[(i)] nine fresh encryptions $\Enc'(D_{ij})$ with decryption
  proofs that $\Enc'(D_{ij}) / \Enc(D_{ij})$ decrypts to $g^0$;
  \item[(ii)] $\Enc(d_{ij})$ with $d_{ij} \in \{0, 1, 2\}$ and the
  disjunction: for $(A, B)$, $\Delta = 2 - 0 = 2$ and the second branch
  $5(2 - 2) + 2 = 2 \in \{1, 2\}$; for $(B, A)$, $\Delta = -2$ and the
  first branch $5 \cdot 0 - 2 = -2 \in \{-2, -1, 0\}$; for $(B, C)$,
  $\Delta = 1$, second branch $5 \cdot 0 + 1 = 1$; for $(C, B)$, first
  branch $-1$;
  \item[(iii)] $\Enc(p_{ij})$ with $p_{ij} - d_{ij} \in \{0, 1, 2\}$ (all
  differences are $0$); the $27$ closure disjunctions, e.g.\ $(B, C)$
  through $A$: $p_{BC} - p_{BA} = 1 \in \{0,1,2\}$ holds (first branch),
  and $(A, C)$ through $B$: $p_{AC} - p_{AB} = 0$ holds; and witness
  paths, here all direct edges, e.g.\ $A \to B$ with $d_{AB} - p_{AB} =
  0$;
  \item[(iv)] for $w = A$: $p_{AB} - p_{BA} = 2$ and $p_{AC} - p_{CA} =
  2$, both in $\{0, 1, 2\}$.
\end{itemize}
A certificate for $B$ would fail at (iv): $p_{BA} - p_{AB} = -2$. A
certificate claiming $p_{CA} = 2$ (to help $C$) would fail at (iii): no
path from $C$ to $A$ has an edge of strength $\ge 2$, since $d_{CA} =
d_{CB} = 0$, so no witness path can be proved.

\section{The tally certificate}\label{sec:tally}

\subsection{What the trustees publish}
Let $N$ be the number of accepted ballots, so $0 \le D_{ij} \le N$. The
trustees decrypt $\Enc(D)$ among themselves, compute $d$, $p$ and the
winner $w$ in private, and publish a certificate over \emph{encrypted}
values. The only primitive it needs is the comparison of two encrypted
integers: $\Enc(u) / \Enc(u')$ encrypts $u - u'$, and $u \ge u'$ is the
membership statement $u - u' \in \{0, \ldots, N\}$, an $(N+1)$-branch
\rocq{EncProof}. The four parts:
\begin{enumerate}[nosep]
  \item[(i)] \emph{Re-encryption.} For every pair, $\Enc'(D_{ij})$, a fresh
  encryption with randomness the trustees know, and a \rocq{DecProof} that
  $\Enc'(D_{ij}) / \Enc(D_{ij})$ decrypts to $g^0$. All later statements
  are about $\Enc'$, whose randomness the trustees know, which is what
  lets them prove things about it.
  \item[(ii)] \emph{Defeats.} $\Enc(d_{ij})$ with a range proof $d_{ij} \in
  \{0, \ldots, N\}$ and a disjunction, with $\Delta = D_{ij} - D_{ji}$ and
  $K = 2N + 1$:
  \[
    \big(K d_{ij} + \Delta \in \{-N, \ldots, 0\}\big) \;\vee\;
    \big(K (d_{ij} - D_{ij}) + \Delta \in \{1, \ldots, N\}\big),
  \]
  the shift trick of Section~\ref{sec:ballot1} again. The first branch is
  possible only for $d_{ij} = 0$ and $\Delta \le 0$; the second only for
  $d_{ij} = D_{ij}$ and $\Delta \ge 1$ (proof below).
  \item[(iii)] \emph{Strongest paths.} $\Enc(p_{ij})$ with: $p_{ij} -
  d_{ij} \in \{0, \ldots, N\}$ for every pair; for every $i, j, k$, the
  disjunction $p_{ij} - p_{ik} \in \{0, \ldots, N\} \vee p_{ij} - p_{kj}
  \in \{0, \ldots, N\}$, i.e.\ $p_{ij} \ge \min(p_{ik}, p_{kj})$; and for
  every pair a witness path $i = u_0, \ldots, u_\ell = j$ with a range
  proof $d_{u_t u_{t+1}} - p_{ij} \in \{0, \ldots, N\}$ on every edge.
  \item[(iv)] \emph{Winner.} $p_{wj} - p_{jw} \in \{0, \ldots, N\}$ for
  every $j \ne w$.
\end{enumerate}
All proofs share one Fiat--Shamir challenge
(\rocq{nizk\_generate\_tally}); the verifier recomputes it and checks every
proof against the aggregate (\rocq{verify\_tally}).

\ex{In the running example, $N = 5$, $K = 11$. Part (ii) for $(A, B)$:
$d_{AB} = 3$, $\Delta = 3 - 2 = 1$; the second branch, $K(3 - 3) + 1 = 1
\in \{1, \ldots, 5\}$, is proved. For $(B, A)$: $d_{BA} = 0$, $\Delta =
-1$; the first branch, $0 - 1 = -1 \in \{-5, \ldots, 0\}$. For $(B, C)$:
$\Delta = 0$, $d = 0$, first branch with $0$. Part (iii): $p_{CB} = 3$,
witness path $C \to A \to B$ with edge proofs $d_{CA} - p_{CB} = 0$ and
$d_{AB} - p_{CB} = 0$; closure disjunction for $(C, B)$ through $A$:
$p_{CB} - p_{CA} = 0 \in \{0, \ldots, 5\}$ (first branch). For $(A, C)$,
$p_{AC} = 0$: witness path $A \to C$ with $d_{AC} - p_{AC} = 0$, and the
disjunction through $B$: $p_{AC} - p_{AB} = -3$ (first branch fails),
$p_{AC} - p_{BC} = 0$ (second branch holds). Part (iv) for $w = C$: $p_{CA}
- p_{AC} = 3$ and $p_{CB} - p_{BC} = 3$. The board holds $9$ re-encryptions,
$9$ defeats, $9$ path strengths, $27$ closure disjunctions, the edge
proofs, and $2$ winner proofs. The certificate for $A$ fails at (iv):
$p_{AC} - p_{CA} = -3$.}

\subsection{The closure lemma}
\begin{lemma}[\rocq{cert\_ok\_is\_closure}, \rocq{closure\_unique}, \rocq{cert\_winner}]
Let $d \ge 0$ and let $p$ and paths satisfy (iii): (1) $d_{ij} \le p_{ij}$;
(2) $\min(p_{ik}, p_{kj}) \le p_{ij}$ for all $i, j, k$; (3) for all $i,
j$, a path from $i$ to $j$ of strength $\ge p_{ij}$ exists. Then $p$ is the
strongest-path matrix of $d$. Consequently the strongest-path matrix is the
unique such $p$, and a $w$ satisfying (iv) is a Schulze winner of $d$.
\end{lemma}
\begin{proof}
Let $p^*$ be the strongest-path matrix. $p \le p^*$: by (3), $p_{ij}$ is at
most the strength of some path, hence at most the maximum $p^*_{ij}$.
$p^* \le p$: we show that every path $i = u_0, \ldots, u_\ell = j$ has
strength $\le p_{ij}$, by induction on $\ell$. For $\ell = 1$ the strength
is $d_{ij} \le p_{ij}$ by (1). For $\ell > 1$, the strength is $\min(d_{i
u_1}, \sigma)$ where $\sigma$ is the strength of the path $u_1, \ldots, j$,
which by induction is $\le p_{u_1 j}$; so the strength is $\le \min(p_{i
u_1}, p_{u_1 j}) \le p_{ij}$ by (1) and (2). Uniqueness: $p^*$ is a
function of $d$. Winner: (iv) says $p_{wj} \ge p_{jw}$ for every $j$, which
is the definition applied to $p = p^*$.
\end{proof}

\subsection{Soundness of the certificate}
\begin{theorem}[\rocq{verify\_tally\_sound}]
Let $h = g^x$, $g$ of order $q$, characteristic above $2N^2 + 3N$, and let
$\Enc(D)$ encrypt a matrix $D$ with $0 \le D_{ij} \le N$. From two
accepting transcripts of a certificate for $w$ on $\Enc(D)$ with the same
announcements and different challenges, $d = \mathrm{defeat}(D)$ and there
are $p$ and paths satisfying (iii) and $w$ satisfying (iv); hence $w$ is a
Schulze winner of $D$.
\end{theorem}
\begin{proof}
Extract every proof as in Theorem~\ref{thm:b1sound}. (i) gives
$\Enc'(D_{ij})$ decrypting to $g^{D_{ij}}$, i.e.\ the re-encryptions carry
$D$. (ii) gives $d_{ij} \in [0, N]$ and one of the two branches. In the
first, $K d_{ij} + \Delta \le 0$ with $\Delta \ge -N$ forces $K d_{ij} \le
N$, so $d_{ij} = 0$ (since $K = 2N + 1 > N$), and then $\Delta \le 0$, so
$D_{ij} \le D_{ji}$ and $\mathrm{defeat}(D)_{ij} = 0 = d_{ij}$. In the
second, $1 \le K(d_{ij} - D_{ij}) + \Delta \le N$ with $|\Delta| \le N$
forces $|K (d_{ij} - D_{ij})| \le 2N < K$, so $d_{ij} = D_{ij}$, and then
$\Delta \ge 1$, so $D_{ij} > D_{ji}$ and $\mathrm{defeat}(D)_{ij} = D_{ij}
= d_{ij}$. The integers compared here have absolute value at most $K N + N
= 2N^2 + 2N$, which the characteristic bound covers, and equality of
plaintexts is read in the exponent using the order of $g$. (iii) and (iv)
give the conditions of the lemma directly, since each is a range statement
on a difference of extracted plaintexts. The lemma concludes.
\end{proof}

\subsection{Zero-knowledge and the end-to-end theorems}
\begin{theorem}[\rocq{tally\_zk}]
The distribution of the certificate over the trustees' randomness is a
permutation of the output of a simulator that has $\Enc(D)$, $w$ and the
challenge only.
\end{theorem}
Every proof of the certificate is a Chaum--Pedersen proof or an
OR-composition of them, simulated as in Section~\ref{sec:sigma}; the
re-encryptions are fresh ElGamal ciphertexts whose hiding is IND-CPA. So
the public learns $w$ and nothing about $D$ beyond what $w$ implies.

\begin{theorem}[\rocq{schulze\_end\_to\_end}, \rocq{schulze\_end\_to\_end2}]
Under the hypotheses of the ballot and certificate soundness theorems,
with $k \le N$ ballots each accompanied by its rewinding transcripts and a
certificate accompanied by its two transcripts on the aggregate of the
ballots: there are position vectors $s^{(1)}, \ldots, s^{(k)}$ (one per
ballot) such that $w$ is a Schulze winner of the margins $D_{ij} = \#\{l :
s^{(l)}_i < s^{(l)}_j\}$; and (for the batched version) this or a binding
break of a commitment key.
\end{theorem}
\begin{proof}
Ballot soundness gives the position vectors; \rocq{cdec\_aggregate} gives
that the aggregate encrypts their margin counts, which are in $[0, k]
\subseteq [0, N]$; certificate soundness on that $D$ gives the winner. The
batched version uses \rocq{ballot2\_sound} and a finite choice with an
escape clause (\rocq{fin\_choice\_or}): either some ballot's extraction
produced a binding break, or every ballot has positions.
\end{proof}

\subsection{Several seats and cost}
Gentoo fills seven council seats: the winner is removed and the method is
run on the remaining candidates. The aggregate restricted to the remaining
candidates is the aggregate of the smaller election, so seven certificates
on the same encrypted margins, of $m, m-1, \ldots, m-6$ candidates, replace
Gentoo's seven passes; the ballots are not touched. The certificate has
$O(m^3)$ proofs, independent of $N$, but each range proof has $N+1$
branches and each closure disjunction $2N + 2$, so its size is $O(m^3 N)$
group elements and the time grows linearly in $N$. On Gentoo's elections
(10--17 candidates, 50--110 voters) the trustees' prover takes from $18$
minutes ($m = 10$, $N = 50$) to $144$ minutes ($m = 17$, $N = 96$) of
processor time and the verifier from $9$ to $71$ minutes, at the
$2048$-bit parameters, single-threaded.

\section{Security, in one place}\label{sec:security}

\paragraph{Verifiability} Ballot verification, aggregation and certificate
verification take only the board and no key. A ballot that fails is posted
as rejected, and the verdict is recomputable. The development packages the
count as an inductive judgement \rocq{count} with rules for a valid ballot,
an invalid ballot and the final certificate, and proves that a derivation
ending in a true verdict, with the rewinding transcripts, names a Schulze
winner of the rankings cast (\rocq{count\_sound}).

\paragraph{Soundness, where the trust lies} The two theorems that decide
the result, ballot soundness and certificate soundness, are about the
extracted verifier. What is \emph{not} covered by a theorem is the wiring
around it: the decoder must check that every group element has order $q$
and every field element is reduced; the verifier must derive the $\delta$
ciphertexts (first ballot) and the challenges $\beta, \gamma, y$ (batched
ballot) itself and never accept them as inputs; and the Fiat--Shamir hash
must bind the voter's credential. The Gentoo driver tests the last by
replaying a ballot under another credential.

\paragraph{Privacy} Ballots are IND-CPA ciphertexts (and perfectly hiding
commitments) with perfectly simulatable proofs; the aggregate is never
decrypted on the board; the certificate is simulatable from $\Enc(D)$ and
$w$. So the public learns $w$, under DDH and an honest threshold of
trustees. The trustees learn $D$ and nothing else: this is \emph{public}
tally-hiding in Kryvos's sense, not full tally-hiding.

\paragraph{Receipt-freeness} The coercer sees the board. Replace the
coerced voter's ranking by any other: the voter's ciphertexts and
commitments are indistinguishable (DDH, perfect hiding), the proofs are
simulated identically, and the certificate is a function of the new
aggregate, which the coercer cannot distinguish from the old unless the
winner changes. So the coercer learns at most what the winner reveals, and
the voter cannot prove a vote: every ranking is consistent with the
voter's board entry. Excluded: a coercer who obtains the device's
encryption randomness, which opens the ciphertexts; that is
coercion-resistance, which needs a re-encryption step by a trusted
authority and is compatible with the scheme.

\section{A recipe for proofs from $\Sigma$-protocol building blocks}\label{sec:recipe}
The two ballots were designed by the same sequence of steps, and the
sequence applies to other problems. This section states it in general,
with the Schulze ballot as the running illustration.

\begin{enumerate}[leftmargin=*]
  \item \textbf{Write the relation precisely.} Separate the public
  statement (keys, ciphertexts, and fixed public maps such as $M$), the
  secret witness (the ranking $s$ and the randomness), and the condition
  that must hold (every encrypted bit equals $[s_i < s_j]$). Most design
  mistakes start with a vague statement of what is being proved.

  \item \textbf{Choose what to encrypt, what to commit, and how to encode
  the secret.} Encrypt only what must later be decrypted or aggregated;
  here that is only the bits $b$, because they are summed into the
  margins. Commit to everything that only needs proving, such as $s$,
  $\delta$, and $p$; commitments are cheaper and support vector arguments.
  Pick an encoding that makes validity easy to state: with a position
  vector every vector is a valid ranking, while a permutation would need a
  shuffle proof. Much of the cost is decided here.

  \item \textbf{Let the verifier derive whatever it can.} Every value the
  verifier computes from public data through the homomorphism is a value
  the prover cannot lie about. In the first ballot the verifier computes
  $\Enc(\delta_{ij}) = \Enc(s_j) / \Enc(s_i)$ itself
  (Section~\ref{sec:ballot1}); in the simplified batched ballot it forms
  $C_0$ from the commitments and uses the derived key $\mathbf{g}'$
  (Section~\ref{sec:ballot2s}). This anchoring is often what rules out the
  attacks.

  \item \textbf{Reduce the condition to the shapes the toolkit proves.}
  The building blocks prove a small set of shapes, listed in
  Table~\ref{tab:shapes}. The work is to rewrite the condition into them,
  and four tricks do most of it: \emph{selection}, writing ``if $b$ then
  $x$ else $y$'' as $b x + (1 - b) y$, which is how $v$ arose;
  \emph{auxiliary witnesses for products}, committing to $p = b \circ
  \delta$ so that everything else stays linear; \emph{ranges via bits},
  turning non-negativity into $v = \sum_t 2^t X^{(t)}$ with each $X^{(t)}$
  a bit vector; and \emph{shifts that merge cases into one range}, as the
  first ballot's $t = K b + \delta$ puts both cases into one membership
  proof.

  \item \textbf{Batch when the same claim repeats.} If a claim must hold
  for many entries, do not prove it entry by entry. A \emph{random linear
  combination} checks one weighted sum with fresh random
  weights~\cite{bgr98}; that is the $\beta$ link. A \emph{polynomial
  identity} encodes $n$ relations as the coefficients of a polynomial in a
  challenge $y$ and checks it at one point~\cite{schwartz,zippel}; that is
  the zero argument. \emph{Vector commitments} let one group element stand
  for a whole vector.

  \item \textbf{Compose and fix the challenge order.} Conjunctions share
  one challenge. Under Fiat--Shamir, hash everything a challenge must not
  be chosen before: the commitments before $\beta$, $\gamma$, and $y$, and
  the announcements before $c$ and $x$. Include the context too, such as
  the voter's credential, which stops replays.

  \item \textbf{Prove the three properties.} Completeness is algebra.
  Special soundness extracts a witness from rewound transcripts (two for a
  $\Sigma$-protocol, $2r + 1$ for an $r$-row zero argument), with the
  conclusion ``a valid witness, or a binding break''. Zero-knowledge needs
  a simulator; in this development it is proved by a bijection on the
  randomness.

  \item \textbf{Check necessity, then optimise.} For every check, find the
  cheat that passes without it, as in Table~\ref{tab:necessary}. A check
  with no such cheat is redundant; a cheat with no check against it is a
  soundness bug. Then count exponentiations and transcript size, and merge
  rows and keys where the homomorphism allows, as in
  Section~\ref{sec:ballot2s}.
\end{enumerate}

\begin{table}[ht]
  \centering
  \caption{The shapes of claim that the building blocks prove.}
  \label{tab:shapes}
  \begin{tabularx}{\textwidth}{@{}>{\raggedright\arraybackslash}X>{\raggedright\arraybackslash}X@{}}
    \toprule
    shape of the claim & building block \\
    \midrule
    knowledge of a preimage under a linear map (``linear in the exponent'') &
    Schnorr~\cite{schnorr}, Okamoto, and Maurer's linear-map
    protocol~\cite{maurer} (Section~\ref{sec:linsigma}) \\
    the same secret behind two values, including correct decryption &
    Chaum--Pedersen~\cite{chaumpedersen} \\
    a value is one of a small list & OR-composition~\cite{cds}
    (Section~\ref{sec:sigma}) \\
    several claims at once & AND: the protocols in parallel with one shared
    challenge \\
    a product of two secrets & a multiplication proof, or a zero argument
    over commitments (Section~\ref{sec:zero}) \\
    a value lies in a range & an OR over the range, or a bit decomposition
    and bit checks \\
    \bottomrule
  \end{tabularx}
\end{table}

The same steps also explain the pitfalls this development ran into, each
of which a reader designing a new proof should check for:
\begin{itemize}[leftmargin=*]
  \item \emph{Missing constraints.} Without the bit check on $b$, a tied
  pair can carry $b = -1$ and subtract one from a margin
  (Table~\ref{tab:necessary}). This is the classic bug of
  under-constrained circuits.
  \item \emph{Wrap-around in the field.} ``$\ge 0$'' means nothing modulo
  $q$, so every integer argument needs a bound; here it is a characteristic
  above $3 \cdot 2^L + 2$ (Theorem~\ref{thm:b2sound}).
  \item \emph{Challenge order.} A $\beta$ known before committing lets a
  voter commit to bits that differ from the encrypted ones.
  \item \emph{Weak Fiat--Shamir.} Leaving the statement or the context out
  of the hash is a known source of real attacks on voting systems.
  \item \emph{Unchecked group membership.} The proofs assume every element
  lies in the prime-order group, so the verifier's decoder must enforce it;
  this is the lesson of the paper's Belenios case study.
  \item \emph{Intuition is weakest where machine-checking helps most.} The
  characteristic bound and the fact that no range check on $s$ is needed
  both came out of the formal proof.
\end{itemize}

\section{Where the ideas come from}\label{sec:sources}
The building blocks are standard. Exponential ElGamal~\cite{elgamal},
Pedersen commitments~\cite{pedersen}, Schnorr's protocol~\cite{schnorr},
the Chaum--Pedersen equality proof~\cite{chaumpedersen}, the
OR-composition of Cramer, Damg{\aa}rd, and Schoenmakers~\cite{cds}, and the
Fiat--Shamir transform~\cite{fiatshamir} are the tools of Helios-style
e-voting, and the Schulze method is Schulze's~\cite{schulze}.

The first ballot's design is new to this work: the position vector as the
encoding of a ranking with ties, the anchoring of every pairwise bit to it
so that the verifier derives $\Enc(\delta_{ij})$ itself, and the shift
trick that turns the consistency disjunction into one membership proof. It
replaces the earlier verifiable Schulze tally of Haines, Pattinson, and
Tiwari~\cite{hpt19}, which checks ballots after a shuffle and decrypts the
margins. The tally certificate descends from the evidence-carrying Schulze
count of Pattinson and Tiwari~\cite{pt17}, whose certificate of the
strongest paths is here proved over encrypted values instead of
published.

The batched ballot combines four known techniques with one new vector.
The library already contained Bayer and Groth's zero
argument~\cite{bayergroth}, mechanised for the Swiss Post shuffle
verifier~\cite{swisspost}; its soundness is the polynomial identity test of
Schwartz and Zippel~\cite{schwartz,zippel}, and its bilinear map is
Groth's~\cite{groth09}. The link from the ciphertexts to the committed bit
vector is Bellare, Garay, and Rabin's batch verification~\cite{bgr98}, as
used by Terelius and Wikstr{\"o}m to tie ElGamal ciphertexts to committed
vectors~\cite{tereliuswikstrom}. The linear argument is Maurer's
$\Sigma$-protocol for the preimage of a linear map~\cite{maurer}, here over
a product of group and field coordinates. The range proof by committed bits
and a bit-ness check is the linear-size form of the range proofs of Bootle
et al.~\cite{bootle16} and Bulletproofs~\cite{bulletproofs}, which compress
the same statement to logarithmic size with an inner-product argument. The
new element is the vector $v = -\delta + 2p - b$, which turns $b_{ij} = [s_i
< s_j]$ into ``$b_{ij}$ is a bit and $v_{ij}$ is a non-negative integer'',
so that validity becomes one linear map, one entry-wise product, and
bit-ness. Kryvos~\cite{kryvos} also proves ballot validity over Pedersen
vector commitments, with Groth16 instead of $\Sigma$-protocols.

\section{What is machine-checked}\label{sec:mech}
Every theorem quoted in this note is proved in the Rocq development and
closed under the global context (no axioms, no admitted lemmas; the audit
script \rocq{bench/assumptions.sh} prints the assumptions of each). The
map:
\begin{center}
\footnotesize
\begin{tabularx}{\textwidth}{@{}l>{\raggedright\arraybackslash}X@{}}
  \toprule
  statement in this note & Rocq file: theorems \\
  \midrule
  linear-map $\Sigma$-protocol, Section~\ref{sec:linsigma} & \rocq{Crypto/LinearSigma.v}: \rocq{lin\_completeness}, \rocq{lin\_soundness}, \rocq{lin\_shvzk} \\
  zero argument, Section~\ref{sec:zero} & \rocq{Crypto/BayerGroth*.v}: \rocq{zero\_completeness}, \rocq{zero\_soundness}, \rocq{zero\_shvzk} \\
  first ballot complete & \rocq{Frontend/Schulze.v}: \rocq{ballot\_valid}, \rocq{nizk\_ballot\_valid} \\
  Theorem~\ref{thm:b1sound} & \rocq{Frontend/Schulze.v}: \rocq{ballot\_sound}; at the instance \rocq{Examples/SchulzeIns.v}: \rocq{ballot\_sound\_ins} \\
  first ballot ZK & \rocq{Frontend/Schulze.v}: \rocq{ballot\_zk} \\
  aggregation & \rocq{Frontend/Schulze.v}: \rocq{aggregate\_margins}, \rocq{cdec\_aggregate} \\
  batched ballot complete & \rocq{Frontend/SchulzeBatched.v}: \rocq{ballot2\_valid}, \rocq{nizk\_ballot2\_valid} \\
  Theorem~\ref{thm:b2sound} & \rocq{Frontend/SchulzeBatched.v}: \rocq{ballot2\_sound}, \rocq{ballot2\_sound\_bits} \\
  batched ballot ZK & \rocq{Frontend/SchulzeBatched.v}: \rocq{ballot2\_zk} \\
  simpler ballot complete, Section~\ref{sec:ballot2sd} & \rocq{Frontend/SchulzeSimple.v}: \rocq{com\_gd}, \rocq{ballot3\_valid}, \rocq{nizk\_ballot3\_valid} \\
  Proposition~\ref{prop:v3sound} & \rocq{Frontend/SchulzeSimple.v}: \rocq{ballot3\_sound}, \rocq{ballot3\_sound\_bits} \\
  simpler ballot ZK & \rocq{Frontend/SchulzeSimple.v}: \rocq{ballot3\_zk} \\
  end to end, simpler ballot & \rocq{Backend/SchulzeSimpleTally.v}: \rocq{schulze\_end\_to\_end3} \\
  closure lemma & \rocq{Backend/SchulzeTally.v}: \rocq{cert\_ok\_is\_closure}, \rocq{closure\_unique}, \rocq{cert\_winner} \\
  certificate complete & \rocq{Backend/SchulzeTally.v}: \rocq{verify\_tally\_complete}, \rocq{nizk\_tally\_complete} \\
  certificate sound & \rocq{Backend/SchulzeTally.v}: \rocq{verify\_tally\_sound} \\
  certificate ZK & \rocq{Backend/SchulzeTally.v}: \rocq{tally\_zk} \\
  end to end, first ballot & \rocq{Backend/SchulzeTally.v}: \rocq{schulze\_end\_to\_end}, \rocq{count\_sound} \\
  end to end, batched ballot & \rocq{Backend/SchulzeBatchedTally.v}: \rocq{schulze\_end\_to\_end2} \\
  side conditions at the instance & \rocq{Utility/SchnorrOrder.v}: \rocq{prime\_order}, \rocq{pow\_injective}, \rocq{natF\_injective\_below\_q} \\
  \bottomrule
\end{tabularx}
\end{center}
What is on paper: the Fiat--Shamir and forking arguments (the mechanised
soundness statements are in the rewinding form), the IND-CPA reduction for
privacy, and the receipt-freeness argument; threshold decryption is not
mechanised (the development uses one trustee key). The honest trustees'
algorithms for the closure and the witness paths are not verified and need
not be: they check their own certificate with \rocq{cert\_ok\_b} before
publishing, and the verifier's soundness does not depend on how the
certificate was produced.

\section{Running it}\label{sec:artefact}
The clients, verifiers and provers are extracted from Rocq to OCaml, with
the Fiat--Shamir hash (SHAKE-256) and the randomness supplied by the
driver. From the repository root, with the default opam switch:
{\small
\begin{verbatim}
dune build
E=_build/default/src/Executable
$E/Schulzecode/main.exe              # the running example, both certificates
$E/SchulzeBenchcode/main.exe 7 100   # first ballot and certificate, m = 7, N = 100
$E/SchulzeBenchcode/batched.exe 17 3 # batched ballot, m = 17, three iterations
python3 bench/gentoo_leak.py --positions   # Gentoo rankings as position vectors
V=bench/gentoo-votes/positions-council-202306.txt
B=bench/gentoo-board/council-202306.txt
$E/SchulzeGentoocode/encrypt.exe $V $B     # the voters: encrypt and prove
$E/SchulzeGentoocode/main.exe $B           # verify, aggregate, certify, verify
\end{verbatim}}
The example program prints the board of the running example and walks the
\rocq{count} derivation ballot by ballot, including the two rejected
ballots of Table~\ref{tab:ballot1} and the rejected certificate for $A$.
The Gentoo programs encrypt a council election's published rankings under
per-voter credentials, verify every ballot, replay one under another
credential (rejected), aggregate, decrypt the margins, certify the winner
and verify the certificate; the winner agrees with Gentoo's on every
election.

\begin{thebibliography}{99}\small
\bibitem{bayergroth} S.~Bayer and J.~Groth. Efficient zero-knowledge argument for correctness of a shuffle. In \emph{EUROCRYPT 2012}, LNCS 7237, pp.~263--280. Springer, 2012.
\bibitem{bgr98} M.~Bellare, J.~A. Garay, and T.~Rabin. Fast batch verification for modular exponentiation and digital signatures. In \emph{EUROCRYPT 1998}, LNCS 1403, pp.~236--250. Springer, 1998.
\bibitem{bootle16} J.~Bootle, A.~Cerulli, P.~Chaidos, J.~Groth, and C.~Petit. Efficient zero-knowledge arguments for arithmetic circuits in the discrete log setting. In \emph{EUROCRYPT 2016}, LNCS 9666, pp.~327--357. Springer, 2016.
\bibitem{bulletproofs} B.~B{\"u}nz, J.~Bootle, D.~Boneh, A.~Poelstra, P.~Wuille, and G.~Maxwell. Bulletproofs: Short proofs for confidential transactions and more. In \emph{IEEE Symposium on Security and Privacy 2018}, pp.~315--334.
\bibitem{chaumpedersen} D.~Chaum and T.~P. Pedersen. Wallet databases with observers. In \emph{CRYPTO 1992}, LNCS 740, pp.~89--105. Springer, 1993.
\bibitem{cds} R.~Cramer, I.~Damg{\aa}rd, and B.~Schoenmakers. Proofs of partial knowledge and simplified design of witness hiding protocols. In \emph{CRYPTO 1994}, LNCS 839, pp.~174--187. Springer, 1994.
\bibitem{elgamal} T.~ElGamal. A public key cryptosystem and a signature scheme based on discrete logarithms. \emph{IEEE Transactions on Information Theory}, 31(4):469--472, 1985.
\bibitem{fiatshamir} A.~Fiat and A.~Shamir. How to prove yourself: practical solutions to identification and signature problems. In \emph{CRYPTO 1986}, LNCS 263, pp.~186--194. Springer, 1987.
\bibitem{groth09} J.~Groth. Linear algebra with sub-linear zero-knowledge arguments. In \emph{CRYPTO 2009}, LNCS 5677, pp.~192--208. Springer, 2009.
\bibitem{hpt19} T.~Haines, D.~Pattinson, and M.~Tiwari. Verifiable homomorphic tallying for the Schulze vote counting scheme. In \emph{VSTTE 2019}, pp.~36--53. Springer, 2019.
\bibitem{kryvos} N.~Huber, R.~K{\"u}sters, T.~Krips, J.~Liedtke, J.~M{\"u}ller, D.~Rausch, P.~Reisert, and A.~Vogt. Kryvos: Publicly tally-hiding verifiable e-voting. In \emph{ACM CCS 2022}, pp.~1443--1457.
\bibitem{maurer} U.~Maurer. Unifying zero-knowledge proofs of knowledge. In \emph{AFRICACRYPT 2009}, LNCS 5580, pp.~272--286. Springer, 2009.
\bibitem{pt17} D.~Pattinson and M.~Tiwari. Schulze voting as evidence carrying computation. In \emph{Interactive Theorem Proving (ITP 2017)}, LNCS 10499, pp.~410--426. Springer, 2017.
\bibitem{pedersen} T.~P. Pedersen. Non-interactive and information-theoretic secure verifiable secret sharing. In \emph{CRYPTO 1991}, LNCS 576, pp.~129--140. Springer, 1992.
\bibitem{schnorr} C.~P. Schnorr. Efficient signature generation by smart cards. \emph{Journal of Cryptology}, 4(3):161--174, 1991.
\bibitem{schulze} M.~Schulze. A new monotonic, clone-independent, reversal symmetric, and Condorcet-consistent single-winner election method. \emph{Social Choice and Welfare}, 36(2):267--303, 2011.
\bibitem{schwartz} J.~T. Schwartz. Fast probabilistic algorithms for verification of polynomial identities. \emph{Journal of the ACM}, 27(4):701--717, 1980.
\bibitem{swisspost} Swiss Post. Swiss Post Voting System: Verifier Specification, version 1.7.1. \url{https://gitlab.com/swisspost-evoting/e-voting/e-voting-documentation}, 2026.
\bibitem{tereliuswikstrom} B.~Terelius and D.~Wikstr{\"o}m. Proofs of restricted shuffles. In \emph{AFRICACRYPT 2010}, LNCS 6055, pp.~100--113. Springer, 2010.
\bibitem{zippel} R.~Zippel. Probabilistic algorithms for sparse polynomials. In \emph{EUROSAM 1979}, LNCS 72, pp.~216--226. Springer, 1979.
\end{thebibliography}

\end{document}
